% concatenated from TCOM.tex
\documentclass[conference, 10pt]{IEEEtran}
%\IEEEoverridecommandlockouts
%% MIGHTY HACK TO FIT TO 5 PAGES (from Y. Polyanskiy)
 \abovedisplayskip=4pt plus 1pt minus 2pt
 \belowdisplayskip=4pt plus 1pt minus 2pt
 \baselineskip=11.00pt
 \setlength{\parskip}{1.5pt plus 1pt minus 2pt}  %backup 4pt plus 1pt minus 2pt
%% MIGHTY HACK TO FIT TO 5 PAGES END
\addtolength{\topmargin}{0.05in} 
\usepackage{setspace}
\usepackage{amsmath}
\usepackage{amssymb}
\usepackage{mathrsfs}
\usepackage{amsthm}
\usepackage{multirow}
\usepackage{color}
\usepackage{array}
\usepackage{booktabs}
\usepackage{float}
\usepackage{url}
\usepackage{comment}
\usepackage{enumerate}
\usepackage{graphicx}
\usepackage{enumitem}
%\usepackage{eucal}
\usepackage{tabularx}
\usepackage{makecell}
%\usepackage{hyperref}
\usepackage[capitalize]{cleveref}
\usepackage{fixfoot}
\usepackage{caption}
\captionsetup[table]{name=Table}
\renewcommand{\thetable}{\arabic{table}}


\usepackage[style=ieee]{biblatex}
\addbibresource{ourbib.bib}

\usepackage{algorithm}
\usepackage[noend]{algorithmic}
\usepackage[table]{xcolor}
\usepackage{blkarray}
\usepackage{balance}

\usepackage{stmaryrd}


\IEEEoverridecommandlockouts


\theoremstyle{plain}
\newtheorem{theorem}{Theorem}
\newtheorem{corollary}[theorem]{Corollary}
\newtheorem{lemma}[theorem]{Lemma}
\newtheorem{proposition}[theorem]{Proposition}

\theoremstyle{definition}
\newtheorem{definition}{Definition}
\newtheorem{example}[definition]{Example}
\newtheorem{problem}{Problem}
\newtheorem{conjecture}[problem]{Conjecture}
\newtheorem{construction}{Construction}
\renewcommand{\theconstruction}{\Alph{construction}}
%\theoremstyle{remark}
\newtheorem{remark}[definition]{Remark}
\newtheorem{note}[definition]{Note}
\newtheorem{case}{Case}

\newcommand{\FF}{\mathbb{F}}
\newcommand{\BB}{\mathbb{B}}
\newcommand{\mCC}{\mathbb{C}}
\newcommand{\bbS}{\mathbb{S}}
\newcommand{\XX}{\mathbb{X}}
\newcommand{\RM}{{\text{\small $\mathbb{RM}$}}}
\newcommand{\HH}{\mathbb{H}}
\newcommand{\DD}{\mathbb{D}}
\newcommand{\sS}{{\sf S}}
\newcommand{\A}{\mathtt{A}}
\newcommand{\T}{\mathtt{T}}
\newcommand{\C}{\mathtt{C}}
\newcommand{\G}{\mathtt{G}}
\newcommand{\GC}{$\mathtt{GC}$}
\newcommand{\onezero}{{\underline{10}\,}}
\newcommand{\zeroone}{{01\,}}
\newcommand{\ins}[1]{{\color{blue}#1}}
 \newcommand\qbin[3]{\textnormal{bin}_{#3}(#1,#2)}
%\newcommand{\cal}{\mathcal}
\newcommand\bbq[1]{\bm{b}_q(#1)}
\newcommand\Bq[1]{\bm{B}_q(#1)}
\newcommand{\numberset}{\mathbb}
\newcommand{\N}{\numberset{N}}
\newcommand{\Z}{\numberset{Z}}
\newcommand{\q}{\numberset{q}}
\newcommand{\R}{\numberset{R}}
\newcommand{\bC}{\mathbb C}
%\newcommand{\mC}{\mathcal{C}}
\newcommand{\K}{\numberset{K}}
\newcommand{\F}{\numberset{F}}
\newcommand{\supp}{\textup{supp}}
\newcommand{\E}{\numberset{E}}
\newcommand{\loc}{\textnormal{loc}}
%\newcommand{\A}{\numberset{A}}
\newcommand{\Ol}{\mathcal{O}}
\newcommand{\mC}{\mathcal{C}}
\newcommand{\mU}{\mathcal{U}}
\newcommand{\mW}{\mathcal{W}}
\newcommand{\mP}{\mathcal{P}}
\newcommand{\mR}{\mathcal{R}}
%\newcommand{\RM}{\textup{RM}}
\newcommand{\mY}{\mathcal{Y}}
\newcommand{\mD}{\mathcal{D}}
\newcommand{\cL}{\mathcal{L}}
\newcommand{\fq}{\F_q}
\newcommand{\fqn}{\F_{q^n}}
\newcommand{\fqm}{\F_{q^m}}
\newcommand{\ps}{\oplus_\textnormal{P}}
\DeclareMathOperator{\diag}{Diag}
\DeclareMathOperator{\wtH}{wt_H}
\DeclareMathOperator{\supH}{sup_H}
\usepackage[style=ieee]{biblatex}
\addbibresource{ourbib.bib}
\DeclareMathAlphabet{\mathbfsl}{OT1}{ppl}{b}{it} %{OT1}{cmr}{bx}{it}

\newcommand{\floor}[1]{{\left\lfloor #1\right\rfloor}}
\newcommand{\ceil}[1]{{\left\lceil #1\right\rceil}}
\DeclareMathOperator*{\argmax}{arg\,max}
\DeclareMathOperator*{\argmin}{arg\,min}

\newcommand{\bbracket}[1]{{\left\llbracket #1\right\rrbracket}}
\newcommand{\equivcyclic}{\underset{\text{cyc}}{\sim}}
\newcommand{\equivrc}{\underset{\text{rc}}{\sim}}

\newcommand{\anina}[1]{{\color{magenta}(Anina: #1)}}
\newcommand{\stas}[1]{{\color{blue}(Stas: #1)}}
\newcommand{\ben}[1]{{\color{red}(Ben: #1)}}
\newcommand{\red}[1]{{\color{red}{#1}}}
\usepackage{balance}

\newcolumntype{Y}{>{\centering\arraybackslash}X} % for centering tables in tabularx
\def\tabularxcolumn#1{m{#1}}
\newcommand{\heading}[1]{\bfseries\begin{tabular}{@{}c@{}} #1 \end{tabular}}


\title{{\huge{Linear Code Conversion in the Merge Regime: General Bounds and Reed--Muller Constructions}
%\title{\vspace{-1ex} {\huge{\textbf{Reducing Coverage Depth in DNA Storage: \\ 
%\vspace{-0.5ex} 
%A Combinatorial Perspective on Random Access Efficiency}}
%Optimizing the Coverage Depth in DNA Storage via 
%The Random Access Coverage Depth Problem
}\vspace{-1ex}}
\author{
    Anina Gruica\IEEEauthorrefmark{1}, 
    Benjamin Jany\IEEEauthorrefmark{2}, 
    and Stanislav Kruglik\IEEEauthorrefmark{3}%
    
    \thanks{\IEEEauthorrefmark{1}Anina Gruica is with the Department of Applied Mathematics and Computer Science, Technical University of Denmark, Kongens Lyngby, Denmark (email: anigr@dtu.dk).}
    \thanks{\IEEEauthorrefmark{2}Benjamin Jany is with the Department of Mathematics and Computer Science, Eindhoven University of Technology, Eindhoven, Netherlands (email: b.jany@tue.nl).}
    \thanks{\IEEEauthorrefmark{3}Stanislav Kruglik is with the Department of Electrical and Photonics Engineering, Technical University of Denmark, Kongens Lyngby, Denmark (email: stakr@dtu.dk).}
    \thanks{Parts of this work have been accepted for presentation at 2026 IEEE International Symposium on Information Theory~\cite{gruica2026convertible}.}\thanks{The work of Anina Gruica was supported by the Villum Fonden under Grant VIL52303 and by the EuroTech Visiting Researcher Program 2025. The work of Benjamin Jany was supported by the Casimir Institute at the Eindhoven University of Technology and by the EuroTech Visiting Researcher Program 2025 and 2026.}
    \thanks{\textit{Corresponding author:} Stanislav Kruglik}
    }




% \title{Convertible Codes for Data and Device Heterogeneity\\[-3mm]}

% \author{
%   \IEEEauthorblockN{
%     Anina Gruica\IEEEauthorrefmark{1},
%     Benjamin Jany\IEEEauthorrefmark{2},
%     Stanislav Kruglik\IEEEauthorrefmark{1}
%   }
%   \IEEEauthorblockA{
%     \IEEEauthorrefmark{1}Technical University of Denmark, Kongens Lyngby, Denmark\\
%     anigr@dtu.dk, stakr@dtu.dk
%   }
%   \IEEEauthorblockA{
%     \IEEEauthorrefmark{2}Eindhoven University of Technology, Eindhoven, Netherlands\\
%     b.jany@tue.nl
%   }
% }

\begin{document}
\date{}

%\addtolength{\topmmargin}{+0.1cm}

\maketitle

%\hspace*{-12pt}

	

	

\begin{abstract}
Erasure codes are a core component of most existing large-scale distributed storage systems, ensuring reliability against node failures. Recent work has shown that adapting code parameters to changing node failure rates can lead to significant storage savings. The default approach is to re-encode the data under a new code, which consumes substantial system resources. Code conversion was introduced to reduce this cost. However, existing work has mainly focused on conversions within specific classes of codes. In this paper, we study scalar linear code conversion in the merge regime for arbitrary linear codes. We derive universal lower bounds on the write and read costs in terms of unchanged and read symbols. The bounds are refined using generalized Hamming weights, which capture support-growth properties of subcodes and can give sharper estimates than minimum-distance-only arguments. We show that the framework recovers known bounds for important special cases and can be strictly stronger when the final code has nontrivial jumps in its generalized Hamming weight hierarchy. We then apply the framework to Reed--Muller codes and construct explicit Reed--Muller convertible codes using the Plotkin decomposition. For a natural Reed--Muller parameter regime, the construction attains the derived write-cost lower bound. For the read cost, the generalized-Hamming-weight analysis is sharp for one initial block, while a gap remains for the other block. 

\end{abstract}
%\textit{A full version~\cite{full} of this paper is accessible at: https://arxiv.org/abs/2301.11730}

\begin{IEEEkeywords}
Codes for Distributed Storage, Algebraic Coding, Code conversion, Reed--Muller codes
\end{IEEEkeywords}

\section{Introduction}

%1. What is the problem?\\
%2. Why is it important?\\
%3. Why is it hard?\\
%4. How is our approach different?\\
%5. Key component of approach?\\
%\ben{Do we need to rewrite the intro and mention our conference paper?}
%Erasure codes are a core component of most existing large-scale distributed storage systems and provide low-overhead tolerance to node failures~\cite{balaji2018erasure, shen2025survey}. In this approach, data are encoded using a linear code with codeword length $n$ and dimension $k$, and the resulting code symbols are distributed across different nodes. The encoding parameters are typically selected based on the current failure rate of the system. However, the failure rates of storage devices may vary over time. Adapting the code rate in response to such variations can lead to significant savings in storage space and, consequently, operational costs~\cite{kadekodi2019cluster}. This property is referred to as {\em device heterogeneity}. Nevertheless, modifying code parameters in a storage system typically requires re-encoding the stored data and, in the case of non-systematic codes, decoding all message symbols. Such procedures involve accessing and reading a large number of stored symbols, which can be costly and may outweigh the benefits of conversion~\cite{maturana2022convertible}.

%To overcome the abovementioned issues, Maturana and Rashmi recently proposed a theoretical code conversion framework that defines the process of re-encoding data from one code to another~\cite{maturana2022convertible}. They also introduced convertible codes that minimize conversion costs while maintaining desired properties, such as the MDS property or locality. Code conversion has been studied in two regimes: the merge regime, in which codewords of the initial code are merged into a smaller number of codewords of the final code, and the opposite split regime. The main focus of this line of work is to optimize the conversion cost either with respect to the number of symbols being read and written during conversion (the so-called access cost~\cite{maturana2022convertible}) or with respect to the amount of information downloaded (the so-called bandwidth cost~\cite{maturana2023bandwidth}). Another line of recent research focuses on security against passive eavesdroppers with access to a subset of stored codeword symbols~\cite{zhang2025secure}. However, the works mentioned above, as well as subsequent papers, primarily focus on two specific families of codes employed in distributed storage systems—codes with locality and MDS codes, along with their generalizations; see~\cite{ge2025locally, ge2024mds, singhvi2025tight, kong2024locally, ramkumar2025mds, shi2025bounds,chopra2024low, krishnan2025linear, wang2025lower} and references therein.

%However, in distributed storage systems, different data objects may attract different levels of user interest and, consequently, may require different target access rates. This property is referred to as {\em data heterogeneity}. At the same time, the total request rate that can be served by each storage node is limited. The problem of characterizing the set of achievable access rates for multiple files under such constraints was recently introduced under the notion of the service rate region~\cite{emina21}. This framework introduces a new perspective on distributed storage system design by characterizing their efficiency and scalability. Recent results connect it with combinatorial design and provide a detailed analysis of the service rate region of Reed–Muller codes, highlighting their suitability for distributed storage systems under heterogeneous request patterns~\cite{ly2025servicelong, ly2025maximal}. At the same time, despite their utility in the service rate region context, Reed–Muller codes have largely remained outside the focus of research on code conversion, and are therefore primarily suited to systems with fixed device characteristics.

%In this paper, we fill this gap by investigating the conversion properties of Reed–Muller codes in the merge regime, offering the first construction of convertible codes that support both data and device heterogeneity. 
%The key elements of our construction are the structural properties of Reed–Muller codes and the Plotkin construction. To evaluate the performance of our conversion procedure, we, for the first time, consider the conversion of general linear codes and derive bounds on the write and read costs by analyzing the matrix representation of the conversion process and by extending puncturing bounds for the conversion of codes with locality from~\cite{ge2025locally, shi2025bounds} to the general case.

%These bounds provide a first-order estimate of the performance of code conversion for general families of codes used in distributed storage systems, as obtaining tighter bounds necessitates exploiting the structural properties of the codes under consideration.
%We furthermore strengthen existing bounds by considering the generalized weights of a code. 
%We emphasize that these bounds constitute a main theoretical contribution of this work, as they establish, for the first time, a unified framework for analyzing arbitrary linear code conversion. The practical relevance of our approach is demonstrated through the analysis of Reed–Muller code conversion, which performs well in the presence of data heterogeneity.

The exponential growth of data stored in large-scale cloud and distributed storage systems, further accelerated by recent AI workloads, has made storage-node failures a routine event rather than an exception. While simple replication provides fault tolerance, it incurs a substantial storage overhead. Erasure coding offers a more storage-efficient alternative and is therefore a core component of modern distributed storage systems. In a typical erasure-coded system, a file is divided into $k$ data symbols, encoded into $n$ coded symbols by means of a linear code, and distributed across $n$ storage nodes. The choice of the code parameters is usually guided by the desired tradeoff between storage overhead, reliability, repair efficiency, and access performance~\cite{balaji2018erasure, shen2025survey}.

Different families of codes optimize different aspects of this tradeoff. Maximum distance separable (MDS) codes provide the largest possible tolerance against node failures for a given redundancy~\cite{roth2006introduction}. Regenerating codes, in comparison, optimize the amount of data transferred during node repair, while codes with locality, or locally repairable codes (LRCs), optimize the number of nodes involved in the repair of a single failed node~\cite{gopalan2012locality}. Reed--Muller codes provide another example of algebraic codes whose relevance to distributed storage systems has recently been highlighted through the
service-rate-region framework, which captures heterogeneous request patterns under
node-capacity constraints~\cite{ly2025servicelong,ly2025maximal}. 
Another family of codes used in storage systems is formed by LRCs with availability, which have multiple repair sets for each symbol~\cite{kruglik2017bounds}, as well as codes with low skip cost~\cite{zhang2025zigzag} and codes with optimal update and access properties~\cite{chen2020enabling}. However, most such constructions are designed for a fixed set of code parameters, corresponding to a fixed target failure rate and a fixed system workload. At the same time, storage-device failure rates and optimization targets may vary over time. It has been shown that adapting the code rate in response to such variations can lead to significant savings in storage space and operational cost~\cite{kadekodi2019cluster}. However, modifying the above-mentioned codes in a distributed storage system typically requires re-encoding the stored data and, in the case of non-systematic codes, decoding all message symbols. Such procedures involve accessing and reading a large number of stored symbols, which can be costly and may outweigh the benefits of conversion~\cite{maturana2022convertible}.


To address this issue, Maturana and Rashmi introduced the framework of \emph{convertible codes}, which formalizes the problem of converting data encoded under one code into data encoded under another code while minimizing the conversion cost~\cite{maturana2022convertible}. Code conversion has mainly been studied in two regimes: the \emph{merge regime}, where several initial codewords are merged into a smaller number of final codewords, and the opposite \emph{split regime}. The efficiency of a conversion procedure is typically measured either by the \emph{access cost}, namely the number of code symbols read and written during conversion~\cite{maturana2022convertible}, or by the \emph{bandwidth cost}, namely the amount of information transferred during conversion~\cite{maturana2023bandwidth}.

A substantial body of work has developed the theory of convertible codes for important code families used in distributed storage. For MDS codes, fundamental lower bounds and optimal constructions are known in the merge regime, with recent work further improving field-size requirements and bandwidth-efficient vector-code constructions~\cite{maturana2020access, maturana2023bandwidth, ramkumar2025mds, chopra2024low}. Code conversion for locally repairable codes has also attracted increasing attention~\cite{shi2025bounds, ge2025locally, kong2024locally}. Recent papers further study generalized merge-convertible codes, improved lower bounds, and optimal constructions for conversion into LRCs. These advances provide a detailed understanding of code conversion for MDS codes, LRCs, and several of their generalizations.

Despite this progress, the existing literature remains largely focused on conversions within specific code families, especially MDS codes and codes with locality. Much less is known about code conversion for arbitrary linear codes. In particular, there is no general framework that gives universal bounds on the access cost of converting between linear codes in the merge regime. Such bounds are important both theoretically and practically: they provide a baseline against which constructions for specific code families can be evaluated, and they identify which parts of the conversion cost are forced by the structure of the initial and final codes. They also cover the largely unexplored setting of conversions between different code families.


%In this paper, we fill this gap by developing a general framework for scalar linear code conversion in the merge regime and applying it to Reed–Muller codes. The main contributions of this paper are as follows. First, we develop a general linear-algebraic framework for scalar linear code conversion in the merge regime. Unlike previous bounds that are tailored to MDS, locally repairable codes, or related structured code families, our bounds apply to arbitrary linear initial and final codes. They are stated in terms of the number of unchanged and read symbols and therefore provide universal lower bounds on the access cost of conversion. Second, we refine these bounds using generalized Hamming weights. This refinement is natural in the conversion problem because unchanged symbols and read symbols impose support constraints on subcodes of the final code. Consequently, generalized Hamming weights capture information that is invisible to minimum-distance-only arguments. We show that the resulting bounds recover known bounds for MDS and LRC codes in special cases and can be strictly stronger when the generalized weight hierarchy of the final code has nontrivial jumps. Third, we proposed an explicit construction of convertible Reed–Muller codes in the merge regime. To the best of our knowledge, Reed–Muller codes have not previously been studied in the context of convertible codes, despite their rich locality structure and their relevance to distributed storage systems with heterogeneous request patterns. Our construction exploits the recursive structure of Reed–Muller codes through the Plotkin construction. We analyze its access cost using the general bounds developed in the first part of the paper. In particular, for a natural parameter regime, the construction attains the derived bounds on the write cost, while the generalized-Hamming-weight bounds provide a sharp estimate of the read cost for one part of the conversion and leave a remaining gap for the other. Thus, the construction demonstrates that Reed–Muller codes can support efficient conversion while preserving their locality advantages, and it identifies read-cost optimality as a concrete open problem.

%The main contribution of this paper is a general access-cost framework for scalar linear code conversion in the merge regime. In contrast to existing lower bounds for MDS convertible codes, locally repairable convertible codes, and related structured families, our bounds do not rely on maximum-distance separability, locality, or a prescribed parity structure. They apply to arbitrary linear initial and final codes and therefore provide a universal baseline for evaluating conversion procedures. A second contribution is the use of generalized Hamming weights in the analysis of conversion cost. This is natural because unchanged coordinates and read coordinates impose support constraints on subcodes of the final code. Minimum-distance arguments only capture the support of one-dimensional subcodes, whereas generalized Hamming weights capture the support growth of higher-dimensional subcodes. As a consequence, the proposed bounds recover known MDS and LRC conversion bounds in important special cases and can be strictly stronger when the final code has nontrivial jumps in its generalized weight hierarchy. A third contribution is an explicit construction of convertible Reed--Muller codes in the merge regime. Reed--Muller codes have recently been connected to distributed storage systems with heterogeneous request patterns, but they have not previously been studied in the framework of convertible codes. The construction exploits the Plotkin decomposition of Reed--Muller codes. In the natural parameter regime of $m=r+2$, it attains the derived write-cost bound. For the read cost, the generalized-Hamming-weight analysis is sharp for the second initial block and improves the bound for the first block, while a remaining gap identifies read-cost optimality as a concrete open problem.

The main contribution of this paper is a universal access-cost framework for scalar linear code conversion in the merge regime. Existing lower bounds for convertible codes are largely tailored to MDS codes, locally repairable codes, or related structured families. In contrast, our bounds apply to arbitrary linear initial and final codes and therefore separate the part of the conversion cost that is forced by linear-algebraic constraints from the part that depends on a specific code family. The main technical point is that unchanged coordinates and read coordinates impose support constraints on subcodes of the final code. Minimum-distance arguments capture only one-dimensional subcodes, whereas generalized Hamming weights capture the support growth of higher-dimensional subcodes. This leads to refined bounds that recover known MDS and LRC conversion bounds in important special cases and can be strictly stronger when the final code has nontrivial jumps in its generalized Hamming weight hierarchy. As an application, we introduce Reed--Muller convertible codes in the merge regime. Reed--Muller codes have recently been connected to distributed storage systems through service-rate and heterogeneous-request models, but they have not previously been studied from the viewpoint of code conversion. Our construction exploits the Plotkin decomposition of Reed--Muller codes and, for a natural parameter regime, attains the derived write-cost lower bound. For the read cost, the generalized-Hamming-weight analysis is sharp for one initial block and gives a substantially stronger lower bound for the other, leaving a concrete read-optimality question for future work. Compared with our conference paper~\cite{gruica2026convertible}, this journal version provides a refined access-cost analysis based on generalized Hamming weights. We formally show that the resulting bounds recover known bounds in important special cases and can be strictly stronger for codes with nontrivial generalized Hamming weight hierarchies. We also expand the comparison between the Reed--Muller conversion construction and the derived bounds, and provide a more detailed treatment of the proofs, preliminaries, and motivation. 




%The contributions are threefold. First, we derive access-cost lower bounds for conversion between arbitrary linear initial and final codes. In contrast to prior bounds for MDS codes, locally repairable codes, and related structured families, our bounds do not rely on maximum-distance-separability, locality assumptions, or a prescribed parity structure, and therefore provide a universal baseline for evaluating conversion procedures for any pair of linear codes.
%Second, we refine these bounds using generalized Hamming weights. This is the natural tool for the conversion problem because unchanged coordinates and read coordinates impose support constraints on subcodes of the final code. Minimum-distance arguments only capture the support of one-dimensional subcodes, whereas generalized Hamming weights measure the support growth of higher-dimensional subcodes. As a result, the proposed bounds recover known MDS and LRC conversion bounds in important special cases and can be strictly stronger when the final code has nontrivial jumps in its generalized weight hierarchy. Third, we construct explicit Reed--Muller convertible codes in the merge regime. Reed--Muller codes have recently been shown to be relevant for distributed storage systems with heterogeneous request patterns, but they have not been studied previously from the viewpoint of code conversion. Our construction exploits the Plotkin decomposition of Reed--Muller codes. In the natural parameter regime of $m=r+2$, the construction attains the derived write-cost bound. For the read cost, the generalized-Hamming-weight analysis is sharp for the second initial block and improves the bound for the first block, while a remaining gap identifies a concrete open problem.

The rest of the paper is organized as follows. Section~\ref{sec2} recalls the required preliminaries on linear codes, convertible codes, and Reed--Muller codes. Section~\ref{sec:gen} develops general lower bounds on the access cost of linear code conversion in the merge regime. Section~\ref{sec:RM} introduces the Reed--Muller conversion construction and compares its performance with the derived bounds. Section~\ref{sec:concl} concludes the paper.

%The rest of the paper is organized as follows. Section~\ref{sec2} provides the required preliminaries. Section~\ref{sec:gen} presents general bounds on the access cost in the merge regime. Convertible Reed–Muller codes are introduced in Section~\ref{sec:RM} and compared with the derived bounds in Section~\ref{sec:comp}. The paper concludes in Section~\ref{sec:concl}.

\section{Convertible Codes in the Merge Regime}\label{sec2}

\subsection{Preliminaries}

In this paper, we let $q$ be a prime power and we denote by $\F_q$ the finite field of $q$ elements. Moreover, unless specified otherwise, we let $k$ and $n$ be integers with $1\le k \le n$. We denote the set $\{1,\dots, n\}$ by $[n]$.

\begin{definition}
    A \textbf{linear code} $\mC$ of \textbf{dimension} $k$ is an $\F_q$-linear subspace of $\F_q^n$ of dimension~$k$. The \textbf{minimum distance} of~$\mC$ is $d(\mC) := \min\{\wtH(x) : x \in \mC, x \ne 0\}$,
    where $\wtH(x):= |\supp(x)|$ is the \textbf{Hamming weight} of $x=(x_1,\dots, x_n) \in \F_q^n$ and $\supp(x) := \{i \in [n] :  x_i \ne 0\}$ its (Hamming) support. We say that a code $\mC$ in $\F_q^n$ of dimension~$k$ is an $[n,k]_q$-code. If additionally $\mC$ has minimum distance $d$, then we say it is an $[n,k,d]_q$-code. Finally we define $\supp(\C) := \bigcup_{c \in \C} \supp(c)$ to be the \textbf{support} of $\C$, and $\C$ is \textbf{non-degenerate} if $\supp(\C) = [n]$. Throughout the paper we assume that codes are non-degenerate.
\end{definition}

%One of the best-known bounds on the paramaters of linear codes states that the length~$n$, the dimension $k$, and the minimum distance $d$, obey the following rule: $k \le n-d+1$. This bound is called \textbf{Singleton bound} and codes meeting this bound with equality are called \textbf{Maximum Distance Seperable} (\textbf{MDS}) codes.


In this paper, we study the \textit{code conversion problem} in the \textit{merge regime} following the theoretical framework introduced in~\cite{maturana2022convertible}, and the generalization thereof from~\cite{ge2024mds}. From now on, we let $\lambda > 1$ be an integer and, without loss of generality, consider the merge of several codes into one. Moreover, we let $1 \le k_F \le n_F$, $k_\textbf{I}:=(k_{I_1},\dots,k_{I_\lambda})$ and $n_\textbf{I}:=(n_{I_1},\dots,n_{I_\lambda})$,  $1 \le k_{I_i} \le n_{I_i}$ for $i \in [\lambda]$ to be integers. We also assume that  $\sum_{i=1}^{\lambda} k_{I_i}=k_F$. We define a \textit{convertible code} formally as follows.

\begin{definition}[\textnormal{see~\cite[Definition 3]{ge2024mds}}]
A \textbf{convertible code} in the merge regime with parameters $(n_\textbf{I},k_\textbf{I},n_F,k_F)$ consists of
    \begin{itemize}
        \item[(i)] $\lambda$ \textbf{initial codes} $\{\mC_{I_i}\}_{i \in [\lambda]}$, where $\mC_{I_i}$ is an $[n_{I_i}, k_{I_i}]_q$ code for all $i \in [\lambda]$;
        \item[(ii)] a \textbf{final code} $\mC_F$ with parameters $[n_F,k_F]_q$;
        \item[(iii)] a \textbf{conversion procedure (map)} $\sigma: \prod_{i=1}^{\lambda} \mC_{I_i} \longrightarrow \mC_F$.  
        %\ben{the following seems a little redundant} that takes as input $(c_{I_1},\dots,c_{I_\lambda}) \in \prod_{i=1}^{\lambda} \mC_{I_i}$ and outputs $\sigma(c_{I_1},\dots,c_{I_\lambda})=c_F \in \mC_F$.
    \end{itemize}
We let $\mC_\textbf{I} :=\prod_{i=1}^{\lambda}\mC_{I_i}$, where elements of $\mC_{\textbf{I}}$ are tuples $(c_1,\dots,c_{\lambda})$ with $c_i \in \mC_{I_i}$ for $i \in [\lambda]$. We denote the convertible code that merges the initial codes $\{\mC_{I_i}\}_{i \in [\lambda]}$ into the final code $\mC_F$ using the conversion procedure $\sigma$, by $(\mC_\textbf{I},\mC_F,\sigma)$. Throughout the paper, unless explicitly stated otherwise, we consider linear information-preserving conversion maps. 
\end{definition}
The main objective in the theory of convertible codes is to merge several \emph{small} codes used in a distributed storage system into a larger code without fully decoding them. There are two main performance metrics, namely access cost and bandwidth. Access cost corresponds to the total number of disks accessed during conversion, while bandwidth reflects the total number of symbols transferred across the network during conversion. To optimize the former, it is sufficient to consider scalar codes, whereas optimizing the latter may require moving to vector codes, since partial downloads can be beneficial for bandwidth savings. In this paper, to sharpen the focus, we consider only the access-cost metric and leave the bandwidth extension as an interesting open problem. 

To better understand the access cost of a conversion procedure, we classify its symbols into three categories. 

\begin{definition}[\textnormal{see~\cite[Section III]{maturana2022convertible}}] \label{def:clasym}
Let $(\mC_\textbf{I},\mC_F,\sigma)$ be a convertible code with parameters $(n_\textbf{I},k_\textbf{I},n_F,k_F)$. Then 
\begin{itemize}
    \item[(i)] for $i \in [\lambda]$, the \textbf{unchanged symbols} corresponding to $\mC_{I_i}$, denoted by $\mU_{i}$, are symbols coming from $\mC_{I_i}$ that remain unchanged under~$\sigma$ (possibly in different locations), and the unchanged symbols from $\mC_{\textbf{I}}$ are $\mU := \bigcup_{i \in [\lambda]} \mU_i$;
    \item[(ii)]  \textbf{new symbols} in $\mC_F$, denoted by $\mW$, are symbols in $\mC_F$ that are not inherited from an unchanged symbol in $\mC_\textbf{I}$;
    %and thus are functions of the read symbols $\cup_{i=1}^{\lambda} \mR_{i}$.
    \item[(iii)] for $i \in [\lambda]$ the \textbf{read symbols} from $\mC_{I_i}$, denoted by $\mR_{i}$, are symbols from $\mC_{I_i}$ that are used to determine a new symbol, and the read symbols from $\mC_{\textbf{I}}$ are $\mR := \bigcup_{i \in [\lambda]}\mR_i$.
    %during the conversion procedure;
    
    
   
\end{itemize}
Clearly, we have the following relation between the number of new and unchanged symbols $|\mW| = n_F - |\mU|$. We note that the sets of unchanged symbols and read symbols are not required to be disjoint. A symbol may be carried over unchanged into the final codeword and also be read in order to compute one or more new symbols. In this case, the symbol is counted both as an unchanged symbol and as a read symbol. The sets of unchanged and read symbols are defined with respect to codeword coordinates, not with respect to a particular generator-matrix representation.
%the number of new symbols of a convertible code is $n^F-|\sum_{i=1}^\lambda \mU_{i}|$.
\end{definition}


We can now formally define the access cost.
\begin{definition}[\textnormal{see~\cite[Definition 5]{maturana2022convertible}}]
Given a convertible code $(\mC_\textbf{I},\mC_F,\sigma)$, the \textbf{read cost} of the conversion procedure counts the number of symbols that are read during the procedure, i.e., $|\mR|$.
%the number of symbols that are used to compute the symbols of the final code. 
The \textbf{write cost} of the conversion procedure counts the number of new symbols in the final code that must be written, i.e., $|\mW|$.
%the number of symbols that are newly computed. 
The \textbf{access cost} of the conversion procedure is the sum of the read and write costs.
\end{definition}

\begin{remark} \label{rem:default}
Let $(\mC_\textbf{I},\mC_F,\sigma)$ be a convertible code with parameters $(n_\textbf{I},k_\textbf{I},n_F,k_F)$. The default conversion procedure reads $k_{I_i}$ code symbols from each initial code $\mC_{I_i}$, respectively, decodes the information, leaves these read symbols unchanged, and writes the new redundancy symbols (there are $n_F-k_F$ of them) to achieve the desired parameters of $\mC_F$. This results in an access cost of $n_F-k_F+\sum_{i=1}^{\lambda}k_{I_i}=n_F-k_F+k_F=n_F$. In this paper, we aim to improve the latter by preserving as many symbols from the initial code as possible and reducing the number of newly generated redundancy symbols during code conversion. 
\end{remark}


In the general framework, no linearity assumption is imposed on the conversion procedure
from $\mathcal C_{\mathbf I}$ to $\mathcal C_F$. However, in order to derive bounds on the
read and write costs, we restrict attention in what follows to linear conversion procedures.
More precisely, we assume that $\sigma:\mathcal C_{\mathbf I}\to \mathcal C_F$ is an $\mathbb F_q$-linear map. We also assume throughout the paper that
$\sum_{i=1}^{\lambda} k_{I_i}=k_F$, so that, under the requirement that the conversion preserves the encoded information, $\sigma$ is an isomorphism between $\mathcal C_{\mathbf I}$ and $\mathcal C_F$. The study of nonlinear conversion procedures is left as an interesting direction for future work.
Equivalently, after choosing generator matrices for the initial and final codes, 
each final
code symbol is a linear combination of the initial code symbols. An unchanged symbol is a
final coordinate that is equal to one of the initial coordinates. A new symbol is a final
coordinate that is not unchanged, and a read symbol is an initial coordinate that is used in
the computation of at least one new symbol.
%to fit 5 page limit
\begin{comment}
We now present an example that illustrates how a carefully designed conversion procedure can reduce the read and write costs relative to the default approach.

\begin{example} \label{ex:conv}
   Let $\mC_{\textbf{I}} : = \mC_{I_1} \times \mC_{I_2} \leq \F_2^8$ and $\mC_F  \leq \F_2^7$, where $\mC_{\textbf{I}}$ and $\mC_F$ have respective generator matrices $G$ and $G_F$ as defined below: $$G_{\textbf{I}}
:= \begin{bmatrix} 1 & 0& 0 & 1 & 0 & 0 &0 &0 \\ 0 & 1 & 0 & 1& 0 & 0 &0 &0 \\ 0 & 0 & 1 &1 & 0 & 0 &0 &0 \\
     0 & 0 &0 &0 &1 & 0& 0 & 1 \\  0 & 0 &0 &0 & 0 & 1 & 0 & 1\\  0 & 0 &0 &0 &0 & 0 & 1 &1  \end{bmatrix}, \quad  \quad G_F := \begin{bmatrix}
       1 & 0 &0 & 0 &0 &0 &1 \\ 
      
      0 & 1 &0 & 0 &0 &0 &1 \\
      0 & 0 & 1 & 0 &0 &0 &1 \\
      0 & 0 &0 & 1 &0 &0 &1 \\
      0 & 0 &0 & 0 &1 &0 &1 \\
      0 & 0 &0 & 0 &0 &1 &1 
   \end{bmatrix}.$$
   An efficient conversion procedure $\sigma : \F_2^8 \rightarrow \F_2^7$ is as follows:
   $$\sigma( (x_1, x_2,x_3,x_4,x_5,x_6,x_7,x_8) ) = (x_1, x_2, x_3, x_5, x_6, x_7, x_4 +x_8).$$
Note that $\sigma$ leaves $6$ symbols unchanged, $x_1, x_2, x_3, x_5, x_6, x_7$ and must write only the last symbol of the final code. Hence the writing cost is $1$. In order to write the last symbol of the final code, the conversion procedure must read $x_4$ and $x_8$, hence the read cost is $2$. This gives a total access cost of $3 < k_F = 6$.
Although those codes may not be used in practice (since $d=2$) this example nicely illustrates how conversion procedures can greatly improve the efficiency of re-encoding data. 
\end{example}
\end{comment}

The following example shows how a carefully designed code conversion can improve the access cost.
% We now present an example that illustrates how a carefully designed conversion procedure can reduce the read and write costs relative to the default approach.
\begin{example} \label{ex:conv}
   Let $\mC_{\mathbf{I}} := \mC_{I_1} \times \mC_{I_2} \subseteq \F_2^6$ and $\mC_F \subseteq \F_2^5$, where $\mC_{\mathbf{I}}$ and $\mC_F$ have respective generator matrices $G_{\mathbf{I}}$ and $G_F$ as defined below:
   \[
   G_{\mathbf{I}}
   :=
   \begin{bmatrix}
      1 & 0 & 1 & 0 & 0 & 0 \\
      0 & 1 & 1 & 0 & 0 & 0 \\
      0 & 0 & 0 & 1 & 0 & 1 \\
      0 & 0 & 0 & 0 & 1 & 1
   \end{bmatrix},
   \;
   G_F :=
   \begin{bmatrix}
      1 & 0 & 0 & 0 & 1 \\
      0 & 1 & 0 & 0 & 1 \\
      0 & 0 & 1 & 0 & 1 \\
      0 & 0 & 0 & 1 & 1
   \end{bmatrix}.
   \]
   An efficient conversion procedure $\sigma : \F_2^6 \rightarrow \F_2^5$ is as follows:
   \[
\sigma((x_1,x_2,x_3,x_4,x_5,x_6)) = (x_1, x_2, x_4,x_5,x_3+x_6).
   \]
   Note that $\sigma$ leaves $4$ symbols, $x_1, x_2, x_4, x_5$, unchanged and must write only the last symbol of the final code. Hence the write cost is $1$. In order to write the last symbol of the final code, the conversion procedure must read $x_3$ and $x_6$, hence the read cost is $2$. This gives a total access cost of $3$.
 By Remark~\ref{rem:default}, the default approach has an access cost of 5.
   % Although those codes may not be used in practice (since $d=2$), this example nicely illustrates how conversion procedures can greatly improve the efficiency of re-encoding data.
\end{example}



%We end this section by making some observations that will be helpful in the later sections.

% We note that a linear conversion procedure can be represented by a conversion matrix, whose properties are formalized below.
% \begin{lemma} \label{lem:matrix}
% Let $(\mC_\textbf{I},\mC_F,\sigma)$ be a convertible code with parameters $(n_\textbf{I},k_\textbf{I},n_F,k_F)$, and for $i \in [\lambda]$ let $G_{I_i} \in \F_q^{k_{I_i} \times n_{I_i}}$ be a generator matrix of $\mC_{I_i}$, respectively, and let $G_{F} \in \F_q^{k_F \times n_F}$ be a generator matrix of $\mC_F$. Consider the block matrix
% \begin{align*}
%     G_{\textbf{I}} = \begin{bmatrix} G_{I_1} & 0 & \cdots & 0 \\ 0 & G_{I_2} & \cdots  & 0 \\
%     \vdots & \vdots && \vdots \\ 0 & 0 & \cdots &G_{I_{\lambda}}\end{bmatrix} \in \F_q^{k_F \times \sum_{i=1}^{\lambda} n_{I_i}}.
% \end{align*}
% The conversion procedure $\sigma$ is linear if and only if it can represented as a matrix $\mY$ with the properties:

% \begin{enumerate}
%     \item[1)] $    G_{\textbf{I}} \cdot \mY = G_F, \quad \mY \in \F_q^{\sum_{i=1}^{\lambda} n_{I_i} \times n_F}.$ 
%     \item[2)] the unchanged symbols of $\sigma$ are in one-to-one correspondence with the supports of columns of $\mY$ of weight 1.
%     \item[3)] the read symbols of $\sigma$ are the initial-code symbols whose coordinates appear in the support of at least one column of $\mY$ of weight at least 2.
% \end{enumerate}
% We call such a matrix $\mathcal{Y}$ a \textbf{conversion matrix} of $(\mC_I, \mC_F, \sigma)$.
% \end{lemma}

% \begin{example}
%     Consider $(\mC_\textbf{I},\mC_F,\sigma)$ and $G_{\mathbf{I}}$ and $G_{F}$ as in \cref{ex:conv}. The conversion matrix with $\mathcal{Y}$ w.r.t. $G_{\mathbf{I}}$ and $G_F$ is given 
%     $$\mathcal{Y} = 
%     \begin{bmatrix} 1 & 0 & 0 & 0 & 0 \\
%                     0 & 1 & 0 & 0 & 0 \\
%                     0 & 0 & 0 & 0 & 1 \\
%                     0 & 0 & 1 & 0 & 0 \\
%                     0 & 0 & 0 & 1 & 0 \\
%                     0 & 0 & 0 & 0 & 1 \\
%     \end{bmatrix}.$$
% From $\mathcal{Y}$ and \cref{lem:matrix} we see that the unchanged symbols are $\{1, 2, 4, 5\}$, and the read symbols are $\{3, 6\}$. 
% \end{example}


\section{General Bounds on Write and Read Cost} \label{sec:gen}
In this section, we fix the initial codes $\mC_{\textbf{I}}=\prod_{i=1}^{\lambda}\mC_{I_i}$, where $\mC_{I_i}$ is an $[n_{I_i},k_{I_i}]_q$-code, and the final $[n_F,k_F]_q$-code $\mC_F$. We further assume for $i\in [\lambda]$ that $\mC_{I_i}$ has minimum distance $d_{I_i}$, the dual of $\mC_{I_i}$, denoted by $\mC_{I_i}^\perp$, has minimum distance $d_{I_i}^\perp$, $\mC_F$ has minimum distance $d_F$, and the dual of $\mC_F$, denoted by $\mC_F^\perp$, has minimum distance $d_F^\perp$.

We recall the definitions of shortening and puncturing of a code, which we will repeatedly use in this paper.
\begin{definition} \label{short/punct}
    Let $\mC \subseteq \F_q^n$ be a code and $S \subseteq [n]$. 
    \begin{itemize}
    %    \item[(i)] $\supp(x):=\{i : x_i \ne 0\}$ denotes the \textbf{support} of $x=(x_1,\dots,x_n) \in \F_q^n$;
        \item[(i)] The \textbf{puncturing} of $\mC$ onto $S$ is $\mC_{|S}:=\{ \pi_S(x): x \in \mC\}$ where $\pi_S:\F_q^n \to \F_q^{|S|}$ is the projection onto the coordinates indexed by $S$.
        \item[(ii)] Let $\mC(S):=\{x \in \mC : \supp(x) \subseteq S\}$. The \textbf{shortening} of $\mC$ by the set $S$ is the code $\pi_S(\mC(S))$.
    \end{itemize}
\end{definition}

Some of our results in this section are inspired by methods from~\cite{ge2025locally, shi2025bounds, kong2024locally}, among others. Unlike those approaches, 
%which are tailored to codes with locality and exploit finer structural information,
our results are more general and rely only on classical parameters such as length, dimension, and (dual) distance.
The main tools in this section are generalized Hamming weights and Wei duality. These allow us to formulate the bounds in terms of standard code parameters while keeping the arguments uniform for arbitrary linear convertible codes. To the best of our knowledge, this is the first use of these tools within the context of convertible codes. This approach can be used to further refine bounds for other classes of convertible codes with known generalized Hamming weight hierarchies, such as codes with locality. However, to sharpen the focus, we consider bounds for general linear codes and leave these extensions as an interesting direction for future research. We will need the following preliminary definitions and results. 

\begin{definition}
Let $\mC\subseteq \F_q^n$ be a linear code of dimension $k$, and let $j\in [k]$. The \textbf{$j$-th generalized Hamming weight} of $\mC$ is $d_j(\mC):=\min\{ |\supp(\mD)| : \mD\le \mC,\ \dim(\mD)=j\}$. 
\end{definition}

We will use two standard facts about generalized Hamming weights.

\begin{lemma}\label{lem:ghw-increasing-rewrite}
Let $\mC\subseteq \F_q^n$ be a linear code of dimension $k$. Then $1\le d_1(\mC)<d_2(\mC)<\cdots<d_k(\mC)\le n$. In particular, for every $j\in [k]$, we have $d_j(\mC)\ge d_1(\mC)+j-1.$
\end{lemma}

\begin{proof}
The strict monotonicity of generalized Hamming weights is standard. The displayed lower bound follows immediately by induction from $d_{t+1}(\mC)\ge d_t(\mC)+1$ for $t\in [k-1]$.
\end{proof}

\begin{theorem}[Wei duality] \cite[Thm. 3]{wei1991generalized}\label{thm:wei-duality-rewrite}
Let $\mC\subseteq \F_q^n$ be an $[n,k]_q$-code. Then $\{d_1(\mC),\dots,d_k(\mC)\}=[n]\setminus \{n+1-d_1(\mC^\perp),\dots,n+1-d_{n-k}(\mC^\perp)\}$. 
\end{theorem}

\subsection{Bounds on the Write Cost}
\label{subsec:boundsunchanged}

In this subsection, we derive upper bounds on the number of unchanged symbols in a convertible code, and hence lower bounds on the write cost. We first bound the individual sets of unchanged symbols $\mU_i$, then derive a bound on the total number of unchanged symbols, and finally compare the obtained bounds.


\begin{proposition}\label{prop:ghw-unchanged-rewrite}
Let $(\mC_{\textbf{I}},\mC_F,\sigma)$ be a convertible code. For $i\in [\lambda]$, let $\mU_i$ be the unchanged symbols corresponding to $\mC_{I_i}$, and define
\begin{align*} k_{-i}:=\sum_{j\in [\lambda]\setminus\{i\}} k_{I_j}=k_F-k_{I_i}. \end{align*}
Moreover, define the subcode
\begin{align*} \widehat{\mC}_i:=\sigma(\mC_{I_1},\dots,\mC_{I_{i-1}},0,\mC_{I_{i+1}},\dots,\mC_{I_\lambda})\subseteq \mC_F. \end{align*}
Then
\begin{align*} |\mU_i|\le \min\{n_{I_i},\, n_F-d_{k_{-i}}(\widehat{\mC}_i)\}. \end{align*}
In particular,
\begin{align*} |\mU_i|\le \min\{n_{I_i},\, n_F-d_{k_{-i}}(\mC_F)\}. \end{align*}
\end{proposition}

\begin{proof}
Without loss of generality, assume $i=1$, and let $\widehat{\mC}:=\sigma(0,\mC_{I_2},\dots,\mC_{I_\lambda})\subseteq \mC_F.$
Clearly, $\dim(\widehat{\mC})=k_{-1}$. Let $N$ be the set of coordinates in the final codeword that are unchanged from $\mC_{I_1}$,
%and hence directly inherited from $\mC_{I_1}$,
or depend only on the read symbols from $\mC_{I_1}$. Let $N^c=[n_F]\setminus N$. If $v\in \widehat{\mC}$, then the first initial code is zero in the definition of $\widehat{\mC}$, so $v|_N=0$. Thus $\supp(\widehat{\mC})\subseteq N^c$, and therefore $d_{k_{-1}}(\widehat{\mC})\le |N^c|=n_F-|N|.$ 
Equivalently,
$|N|\le n_F-d_{k_{-1}}(\widehat{\mC})$.
Since $\mU_1\subseteq N$, we obtain
$|\mU_1|\le |N|\le n_F-d_{k_{-1}}(\widehat{\mC}).$
Together with the trivial bound $|\mU_1|\le n_{I_1}$, this proves the first claim.

For the second claim, note that $\widehat{\mC}_i$ is a $(k_{-i})$-dimensional subcode of $\mC_F$, and hence
$d_{k_{-i}}(\mC_F)\le d_{k_{-i}}(\widehat{\mC}_i).$
Therefore,
$n_F-d_{k_{-i}}(\widehat{\mC}_i)\le n_F-d_{k_{-i}}(\mC_F),$
which proves the stated weaker bound depending only on $\mC_F$.
\end{proof}

\begin{corollary}\label{cor:unchanged1}
Let $(\mC_{\textbf{I}},\mC_F,\sigma)$ be a convertible code. For $i\in [\lambda]$, let $\mU_i$ be the unchanged symbols from $\mC_{I_i}$. Then
\begin{align*} |\mU_i|\le \min\{n_{I_i},\, n_F-d_F-\textstyle\sum_{j\ne i} k_{I_j}+1\}. \end{align*}
\end{corollary}

\begin{proof}
By Proposition~\ref{prop:ghw-unchanged-rewrite} we have
\begin{align*} |\mU_i|\le \min\{n_{I_i},\, n_F-d_{k_{-i}}(\mC_F)\}. \end{align*}
By Lemma~\ref{lem:ghw-increasing-rewrite},
\begin{align*} d_{k_{-i}}(\mC_F)\ge d_1(\mC_F)+k_{-i}-1=d_F+k_{-i}-1. \end{align*}
Therefore we have
\begin{align*}
|\mU_i|&\le \min\{n_{I_i},\, n_F-(d_F+k_{-i}-1)\}\\&
=\min\{n_{I_i},\, n_F-d_F-k_{-i}+1\} \\
&=\min\{n_{I_i},\, n_F-d_F-\textstyle\sum_{j\ne i} k_{I_j}+1\},
\end{align*}
which proves the corollary.
\end{proof}


\begin{remark}\label{rem:write-ghw-improvement}
Fix $i\in [\lambda]$ and write $k_{-i}=\sum_{j\ne i}k_{I_j}$. Proposition~\ref{prop:ghw-unchanged-rewrite} gives
\begin{align*}
    |\mU_i|\le \min\{n_{I_i},\,n_F-d_{k_{-i}}(\mC_F)\},
\end{align*}
whereas Corollary~\ref{cor:unchanged1} gives the distance-based bound
\begin{align*}
    |\mU_i|\le \min\{n_{I_i},\,n_F-d_F-k_{-i}+1\}.
\end{align*}
Since generalized Hamming weights satisfy $d_{k_{-i}}(\mC_F)\ge d_F+k_{-i}-1$, the bound taking into account generalized Hamming weights is always at least as strong. Moreover (ignoring the minimum with $n_{I_i}$) it is strictly stronger if and only if
\begin{align*}
    d_{k_{-i}}(\mC_F)>d_F+k_{-i}-1.
\end{align*}
Equivalently, the generalized Hamming weights of $\mC_F$ grow faster than the minimal growth forced by strict monotonicity between the first and the $(k_{-i})$-th generalized Hamming weight. In particular, if one of the increments
\begin{align*}
    d_{t+1}(\mC_F)-d_t(\mC_F),
\end{align*}
for $1\le t<k_{-i}$, is greater than $1$, then the bound taking into account generalized Hamming weights gives an improvement over the bound obtained from the minimum distance alone.
\end{remark}

\begin{remark}\label{remark:unchanged-LRC}
Let us compare Proposition~\ref{prop:ghw-unchanged-rewrite} with the improved bound for
locally repairable convertible codes from~\cite{ge2025locally} in the special case of
optimal locally repairable codes. In this setting, closed-form estimates can be obtained
from known lower bounds on generalized Hamming weights~\cite{hao2020generalized}. For comparison, we use the following Singleton-type bound on the minimum
distance of codes with locality $r$ from~\cite{gopalan2012locality}: $$d \leq n-k+1-\left(\left\lceil \frac{k}{r}\right\rceil-1\right).$$ Assume
%, without loss of generality,
that $k_{I_i}=k$ for all \(i\in[\lambda]\), \(n_F=n\), $k_F=\lambda k$ and that all involved codes
have locality $r$. In this case, the bound from \cite{ge2025locally} takes the following form: $$|\mU_i|
 \leq
 k+\left\lceil \frac{\lambda k}{r}\right\rceil
 -\left\lceil \frac{(\lambda-1)k}{r}\right\rceil.$$ Employing the following lower bound on the
$j$-th generalized Hamming weight of an optimal code with locality $r$, length $n$ and dimension $k$ from~\cite[Theorem~4]{hao2020generalized}: %\ben{I don't get the following lower bound. where does it come from. If $\mC$ is optimal LRC then $d = n-k+1-(\lceil \frac{k}{r} \rceil - 1)$ so then shouldn't it be 
%$$d_j \geq d + j -1 = n-k - \lceil \frac{k}{r} \rceil +1 +j $$}
$$d_j
\geq
n-k
-
\left\lceil
\frac{\left\lceil k/r\right\rceil r-j+1}{r}
\right\rceil
+j+1$$ our bounds take the following shape: $$|\mU_i|\leq
\left\lceil
\frac{
\left\lceil \lambda k/r\right\rceil r-(\lambda-1)k+1
}{r}
\right\rceil
+k-1.$$ Subtracting one from another,
and denoting $A=\left\lceil \frac{\lambda k}{r}\right\rceil$, $B=\frac{(\lambda-1)k}{r}$, we get $A-\lceil B\rceil
-
\lceil A-B+\frac{1}{r}\rceil
+1$. Since $A$ is an integer, we have $A-\lceil B\rceil=\lfloor A-B\rfloor$. Therefore the difference becomes $\left\lfloor A-B\right\rfloor + 1
-
\lceil A-B+\frac{1}{r}\rceil$. If $0 < \lceil A - B \rceil - (A-B) < 1/r$ then the expression is $-1$, otherwise the expression is zero.
However, defining $M=Ar-(\lambda-1)k$, we can show that $\lceil A - B \rceil - (A-B)=\lceil
\frac{M}{r}\rceil-\frac{M}{r}$ cannot lie in the open interval from $0$ to $1/r$, since $M$ is an integer. As a result, the expression becomes zero and, in this special case, our generalized-Hamming-weight bound is always as good as the bound from~\cite{ge2025locally}. 

However, the above comparison uses only a lower bound on the generalized Hamming
weights. Therefore, our estimate is sharp only when this lower bound is tight. If the final code has larger jumps in its generalized Hamming weight hierarchy than those guaranteed
by the general lower bound, then Proposition~\ref{prop:ghw-unchanged-rewrite} can be strictly stronger.
For example, consider the binary simplex code of dimension $3$ and length $7$. This code
has locality $2$. In this case, our bound gives $|\mU_i|\leq 1$, whereas the bound from~\cite{ge2025locally} gives only $|\mU_i|\leq 2$.
\end{remark}


A sharper (but more restrictive) bound is as follows. It also takes into account the dual distance of the final code.

\begin{corollary}\label{cor:unchanged2}
Let $(\mC_{\textbf{I}},\mC_F,\sigma)$ be a convertible code. Fix $i\in [\lambda]$, and let $\mU_i$ be the unchanged symbols from $\mC_{I_i}$. If $d_F^\perp>k_{I_i}+1$, then
\begin{align*} |\mU_i|\le k_{I_i}. \end{align*}
In particular, if $d_F^\perp>k_{I_i}+1$ for all $i\in [\lambda]$, then
\begin{align*} |\mU|\le \sum_{i=1}^{\lambda} k_{I_i}=k_F. \end{align*}
\end{corollary}

\begin{proof}
By Proposition~\ref{prop:ghw-unchanged-rewrite}, it suffices to show that $d_{k_{-i}}(\mC_F)\ge n_F-k_{I_i}.$
Since $d_F^\perp=d_1(\mC_F^\perp)>k_{I_i}+1$, we have for every $s\in [n_F-k_F]$,
\begin{align*} d_s(\mC_F^\perp)\ge d_1(\mC_F^\perp)>k_{I_i}+1. \end{align*}
Hence we have
$n_F+1-d_s(\mC_F^\perp) < n_F-k_{I_i}$ for every such $s$ and $\{n_F-k_{I_i},\, n_F-k_{I_i}+1,\,\dots,\, n_F\} \cap  \{n_F+1-d_1(\mC_F^\perp),\dots,n_F+1-d_{n_F-k_F}(\mC_F^\perp)\} = \emptyset$. Therefore, by Wei duality,  $\{n_F-k_{I_i},\, n_F-k_{I_i}+1,\,\dots,\, n_F\}  \subseteq \{d_1(\mC_F),\dots,d_{k_F}(\mC_F)\}.$ 
There are exactly $k_{I_i}+1$ such integers. Since $k_{-i}=k_F-k_{I_i}$, it follows that the $(k_F-k_{I_i})$-th generalized Hamming weight satisfies
$d_{k_{-i}}(\mC_F)\ge n_F-k_{I_i}.$ 
Consequently, $|\mU_i|\le n_F-d_{k_{-i}}(\mC_F)\le k_{I_i}.$
Summing over $i$ proves the final claim.
\end{proof}

%Note that Corollary~\ref{cor:unchanged2} can be proven also directly by considering the dual distance of the final code:     First note that any $t \leq d^{\perp}_F-1$ columns of a generator matrix of $\C_F$ must be linearly independent. Let $i \in [\lambda]$. If $|\mU_{I_i}| \geq k_{I_i}+1$, then because $\dim(\C_{I_i}) = k_{I_i}$ it must be that the columns indexed by elements of $|\mU_i|$ in $\C_F$ are linearly dependent. However, that would imply that $d_F \leq k_{I_i} +1$, a contradiction.

The bound from Corollary~\ref{cor:unchanged2} applied to the case where the final code is MDS recovers~\cite[Lemma 11]{maturana2022convertible}.

\begin{corollary}\label{cor:finalMDSunchangup}
    Let $(\C_{\textbf{I}}, \C_F, \sigma)$ be a convertible code in which $\C_F$ is MDS. For $i \in [\lambda]$, let~$\mU_i$ be the unchanged symbols of $\C_{I_i}$. Then $\sum_{i=1}^{\lambda}|\mU_i| \leq k_F.$
\end{corollary}
\begin{proof}
If $\C_F$ is MDS, then its dual $\mC_F^\perp$ is MDS as well, and so $d^{\perp}_F =n_F-(n_F-k_F)+1= k_F + 1 > k_{I_i} + 1$ for all $i \in [\lambda]$. Hence we can apply Corollary~\ref{cor:unchanged2} to obtain the statement of the corollary.
\end{proof}

In Corollary~\ref{cor:unchanged2}, it is not possible to omit the requirement that $d^{\perp}_F > k_{I_i} + 1$ to derive a non-trivial bound on the number of unchanged coordinates without involving additional parameters of the code. This is illustrated in the following example. 

\begin{example}   
 Let $\lambda=2$ and let
 \begin{align*}
 &\C_{I_1}= \C_{I_2} = \left\langle \begin{bmatrix} 1 & 1 \end{bmatrix} \right\rangle, \quad \C_{\textbf{I}}  = \C_{I_1} \times \C_{I_2},\\
 &\C_F = \left\langle \begin{bmatrix} 1 & 0 & 1 \end{bmatrix}, \begin{bmatrix} 0 & 1 & 0  \end{bmatrix}  \right\rangle.
 \end{align*}
Note that $d_F^{\perp} = 2 = k_{I_i} +1$ for $i \in [2]$. The code $\C_F$ can be obtained by leaving the two coordinates of the first factor and the first coordinate of the second factor unchanged.
\end{example}

    
Let us also derive a bound on the total number of unchanged symbols.

\begin{proposition}\label{prop:totalunchanged}
Let $(\mC_{\textbf{I}},\mC_F,\sigma)$ be a convertible code with $\lambda\ge 2$. Let $\mU$ be the total set of unchanged symbols from $\mC_{\textbf{I}}$. Then
\begin{align*} |\mU|\le \frac{\lambda n_F-k_F-\lambda(d_F-1)}{\lambda-1}. \end{align*}
\end{proposition}

\begin{proof}
Let $i\in [\lambda]$, and consider the subcode $\widetilde{\mC}_i:=\sigma(0,\dots,0,\mC_{I_i},0,\dots,0)\subseteq \mC_F$. Note that $\supp(\widetilde{\mC}_i)\subseteq \mU_i\cup \mW$ and $\dim(\widetilde{\mC}_i)=k_{I_i}$. Let $S=\mU_i\cup \mW$. Clearly, $d(\widetilde{\mC}_i)\ge d_F$, so applying the Singleton bound to $\widetilde{\mC}_i$ gives
\begin{align*} d_F\le d(\widetilde{\mC}_i)\le |S|-\dim(\widetilde{\mC}_i)+1=|S|-k_{I_i}+1. \end{align*}
Since $|S|=|\mU_i|+|\mW|$, summing over $i\in [\lambda]$ gives
\begin{align*}
\sum_{i=1}^{\lambda}|\mU_i|+\lambda|\mW|&\ge \sum_{i=1}^{\lambda}k_{I_i}+\lambda(d_F-1), \\
|\mU|+\lambda|\mW|&\ge k_F+\lambda(d_F-1).
\end{align*}
Now $|\mW|=n_F-|\mU|$, so $|\mU|+\lambda(n_F-|\mU|)=\lambda n_F+(1-\lambda)|\mU|\ge k_F+\lambda(d_F-1)$, which yields the claimed upper bound.
\end{proof}

We now compare the upper bounds on the number of unchanged symbols obtained above.

\begin{remark}\label{rem:ghw-write-comparison-rewrite}
Summing the bounds from Corollary~\ref{cor:unchanged1} over all $i\in [\lambda]$ gives $|\mU|\le \lambda n_F-\lambda d_F-(\lambda-1)k_F+\lambda$. 
Comparing this with Proposition~\ref{prop:totalunchanged}, we obtain
\begin{align*}
&\left(\lambda n_F-\lambda d_F-(\lambda-1)k_F+\lambda\right)-\frac{\lambda n_F-k_F-\lambda(d_F-1)}{\lambda-1}\\ &=\frac{\lambda(\lambda-2)}{\lambda-1}(n_F-k_F-d_F+1)\ge 0,
\end{align*}

where the last inequality follows from the Singleton bound for $\mC_F$. Hence Proposition~\ref{prop:totalunchanged} is always at least as strong as the global bound obtained by summing the bounds from Corollary~\ref{cor:unchanged1}. 

Summing the bound from Proposition~\ref{prop:ghw-unchanged-rewrite} over all $i\in[\lambda]$ gives 
\begin{align*} 
&|\mU| \leq \sum_{i=1}^{\lambda}\left(n_F-d_{k_{-i}}(\mC_F)\right)\\
&= \lambda n_F-\lambda d_F-(\lambda-1)k_F+\lambda -\sum_{i=1}^{\lambda}\Delta_i, \end{align*}
where $\Delta_i = d_{k_-i}(\mathcal C_F)-\bigl(d_F+k_{-i}-1\bigr) \geq 0$. 

Let $\rho_F=n_F-k_F-d_F+1$ be the Singleton defect of $\mC_F$. Comparing the above bound with the bound from Proposition~\ref{prop:totalunchanged}, we obtain 
\begin{align*} &\frac{\lambda n_F-k_F-\lambda(d_F-1)}{\lambda-1} \\
&-\left( \lambda n_F-\lambda d_F-(\lambda-1)k_F+\lambda -\sum_{i=1}^{\lambda}\Delta_i \right) \notag\\ & = \sum_{i=1}^{\lambda}\Delta_i -\frac{\lambda(\lambda-2)}{\lambda-1} \bigl(n_F-k_F-d_F+1\bigr)\\
&= \sum_{i=1}^{\lambda}\Delta_i -\frac{\lambda(\lambda-2)}{\lambda-1}\rho_F . \end{align*}

Hence, the bound from Proposition~\ref{prop:ghw-unchanged-rewrite} is at least as strong as the bound from Proposition~\ref{prop:totalunchanged} whenever $\sum_{i=1}^{\lambda}\Delta_i \geq \frac{\lambda(\lambda-2)}{\lambda-1}\rho_F$.  In particular, if $\lambda=2$, then the right-hand side is zero. Therefore, for $\lambda=2$, the bound from Proposition~\ref{prop:ghw-unchanged-rewrite} is always at least as strong as the bound from Proposition~\ref{prop:totalunchanged}. For $\lambda>2$, the situation depends on the jumps in the generalized Hamming weight hierarchy. If these jumps are sufficiently large, then the bound from Proposition~\ref{prop:ghw-unchanged-rewrite} can be stronger than the unconditional bound from Proposition~\ref{prop:totalunchanged}. For MDS codes, all these bounds coincide.

We next compare the bounds from Proposition~\ref{prop:ghw-unchanged-rewrite}, Corollary~\ref{cor:unchanged1} and Proposition~\ref{prop:totalunchanged} with the more restrictive bound from Corollary~\ref{cor:unchanged2}. The latter bound states that if $d_F^\perp>k_{I_i}+1$ for all $i\in [\lambda]$, then $|\mU|\le k_F$. It is sharper than the bound from Proposition~\ref{prop:totalunchanged} and, consequently, also sharper than
the bound obtained by summing Corollary~\ref{cor:unchanged1}, since
$\frac{\lambda n_F-k_F-\lambda(d_F-1)}{\lambda-1}-k_F=\frac{\lambda}{\lambda-1}(n_F-k_F-d_F+1)\ge 0$, 
again by the Singleton bound. 
We now compare the bound from Proposition~\ref{prop:ghw-unchanged-rewrite} with the bound obtained from Corollary~\ref{cor:unchanged2} that, as we saw before, can improve on the bound from Proposition~\ref{prop:totalunchanged}.

\begin{align*} &\left( \lambda n_F-\lambda d_F-(\lambda-1)k_F+\lambda -\sum_{i=1}^{\lambda}\Delta_i \right)-k_F\\ &= \lambda\bigl(n_F-k_F-d_F+1\bigr) -\sum_{i=1}^{\lambda}\Delta_i= \lambda\rho_F-\sum_{i=1}^{\lambda}\Delta_i . \end{align*} 

As a result, when the generalized Hamming weights have sufficiently large jumps, the bound from Proposition~\ref{prop:ghw-unchanged-rewrite} is at least as strong as all the other bounds. If this is not the case, the bound from Corollary~\ref{cor:unchanged2} can be the best one whenever the additional dual-distance assumptions are satisfied, while the bound from Proposition~\ref{prop:totalunchanged} can give the strongest unconditional upper bound.
\end{remark}

\subsection{Bounds on the Read Cost}
\label{subsec:boundsread}

In this subsection, we derive lower bounds on the number of read symbols. We first prove a generalized-Hamming-weight bound and then obtain two useful special cases as corollaries. Before proceeding further, for the reader’s convenience, we prove a generalized Hamming weight version of the well-known puncturing lemma. 

\begin{lemma}\label{lem:puncturing-ghw-rewrite}
Let $\mC\subseteq \F_q^n$ be a code of dimension $k$, and let $j\in [k]$. If $S\subseteq [n]$ satisfies $|S|\ge n-d_j(\mC)+1$, then the punctured code $\mC|_S$ has dimension at least $k-j+1$.
\end{lemma}

\begin{proof}
Let $T:=[n]\setminus S$, and consider the puncturing map $\pi_S\colon \mC\longrightarrow \mC|_S.$
Its kernel is $\mC(T):=\{x\in \mC : \supp(x)\subseteq T\}.$
Hence, by rank--nullity,
$\dim(\mC|_S)=k-\dim(\mC(T))$. 

Assume for contradiction that $\dim(\mC|_S)\le k-j$. Then $\dim(\mC(T))\ge j$, so $\mC(T)$ contains a $j$-dimensional subcode $\mD\le \mC$. Since $\supp(\mD)\subseteq T$, we obtain
\begin{align*} d_j(\mC)\le |\supp(\mD)|\le |T|=n-|S|. \end{align*}
On the other hand, the assumption on $S$ implies
\begin{align*} n-|S|\le d_j(\mC)-1. \end{align*}
This is a contradiction. Therefore, $\dim(\mC|_S)\ge k-j+1$.
\end{proof}


\begin{proposition}\label{prop:readlower}
Let $(\mC_{\textbf{I}},\mC_F,\sigma)$ be a convertible code. Fix $i\in [\lambda]$, and let $\mR_i$ and $\mU_i$ be the sets of read and unchanged symbols from $\mC_{I_i}$, respectively. Then
\begin{align*} |\mR_i|\ge \max_{1\le j\le k_{I_i}}\left\{ k_{I_i}-j+1-\bigl(|\mU_i|-d_j(\mC_F)+1\bigr)^+ \right\}, \end{align*}
where $(a)^+:=\max\{a,0\}$.
\end{proposition}

\begin{proof}
Without loss of generality, assume $i=1$. Consider the subcode $\widetilde{\mC}:=\sigma(\mC_{I_1},0,\dots,0)\subseteq \mC_F$. Clearly, $\dim(\widetilde{\mC})=k_{I_1}$. Moreover, since $\widetilde{\mC}\le \mC_F$, for every $j\in [k_{I_1}]$ we have
$d_j(\widetilde{\mC})\ge d_j(\mC_F).$

Fix $j\in [k_{I_1}]$, and choose a subset $T\subseteq \mU_1$ of size $|T|=\min\{|\mU_1|,\, d_j(\mC_F)-1\}$. 
Let $S:=[n_F]\setminus T$. Then $|S|=n_F-|T|=n_F-\min\{|\mU_1|,\, d_j(\mC_F)-1\}\ge n_F-d_j(\mC_F)+1\ge n_F-d_j(\widetilde{\mC})+1$. Applying Lemma~\ref{lem:puncturing-ghw-rewrite} to $\widetilde{\mC}$, we obtain
$\dim(\widetilde{\mC}|_S)\ge k_{I_1}-j+1$. 

Now, the nonzero coordinates of $\widetilde{\mC}|_S$ are determined either by the read symbols in $\mR_1$ or by the unchanged symbols in $\mU_1\setminus T$. Hence $\dim(\widetilde{\mC}|_S)\le |\mR_1|+|\mU_1\setminus T|$. Combining the two inequalities gives $|\mR_1|\ge k_{I_1}-j+1-|\mU_1\setminus T|$. Finally,
$|\mU_1\setminus T|=|\mU_1|-\min\{|\mU_1|,\, d_j(\mC_F)-1\}=\bigl(|\mU_1|-d_j(\mC_F)+1\bigr)^+$. Therefore we have
$|\mR_1|\ge k_{I_1}-j+1-\bigl(|\mU_1|-d_j(\mC_F)+1\bigr)^+$. Since this holds for every $j\in [k_{I_1}]$, taking the maximum over $j$ completes the proof.
\end{proof}

\begin{theorem}\label{thm:readlowergen}

Let $(\mC_{\mathbf I},\mC_F,\sigma)$ be a convertible code. Fix $i\in[\lambda]$, and let $\mR_i$ and $\mU_i$ be the sets of read and unchanged symbols from $\mC_{I_i}$, respectively. If there exists an index $j\in[k_{I_i}]$ such that $d_j(\mC_F)<|\mU_i|+1$, let $j$ be the largest such index. Then \[ |\mR_i| \geq \begin{cases} k_{I_i}-j, & \text{if such an index } j \text{ exists},\\ k_{I_i}, & \text{otherwise}. \end{cases} \]
\end{theorem}


\begin{proof}
If no index $j\in[k_{I_i}]$ satisfies $d_j(\mC_F)<|\mU_i|+1$, then
$d_\ell(\mC_F)\ge |\mU_i|+1$ for all $\ell\in[k_{I_i}]$. Hence $(|\mU_i|-d_\ell(\mC_F)+1)^+=0$ for all $\ell$, and the maximum in Proposition~\ref{prop:readlower} is achieved at $\ell=1$, with value $k_{I_i}$. 

Now assume that such an index exists, and let $j\in[k_{I_i}]$ be the largest index such that
$d_j(\mC_F)<|\mU_i|+1$. If $j=k_{I_i}$, then the desired bound is
$|\mR_i|\ge k_{I_i}-j=0$, which is trivial. Thus, in the rest of the proof,
we may assume that $j<k_{I_i}$.

We compare the terms in the maximum from \cref{prop:readlower}.
For $\ell\le j$, we have
$|\mU_i|-d_\ell(\mC_F)+1\ge |\mU_i|-d_j(\mC_F)+1>0$. Moreover, by strict
monotonicity of the generalized Hamming weights,
$d_j(\mC_F)-d_\ell(\mC_F)\ge j-\ell$. Hence
\begin{align*}
&k_{I_i}-j+1-\bigl(|\mU_i|-d_j(\mC_F)+1\bigr)^+\\
&=
k_{I_i}-j+1-|\mU_i|+d_j(\mC_F)-1 \\
&\ge
k_{I_i}-j+1-|\mU_i|+j-\ell+d_\ell(\mC_F)-1 \\
&=
k_{I_i}-\ell+1-\bigl(|\mU_i|-d_\ell(\mC_F)+1\bigr)^+ .
\end{align*}
Thus, among all $\ell\le j$, the largest possible value is bounded above by the term with $\ell=j$.

Since $j$ is maximal and $j<k_{I_i}$, we have
$d_\ell(\mC_F)\ge |\mU_i|+1$ for all $\ell>j$. Therefore
$(|\mU_i|-d_\ell(\mC_F)+1)^+=0$ for all $\ell>j$, and hence $k_{I_i}-\ell+1-\bigl(|\mU_i|-d_\ell(\mC_F)+1\bigr)^+
=
k_{I_i}-\ell+1
\le
k_{I_i}-j$ for all $\ell>j$. In particular, for $\ell=j+1$, we obtain $k_{I_i}-(j+1)+1-\bigl(|\mU_i|-d_{j+1}(\mC_F)+1\bigr)^+
=
k_{I_i}-j$.

It remains to compare this value with the maximum over $\ell\le j$. Since
$d_j(\mC_F)<|\mU_i|+1$ and all quantities are integers, we have
$(|\mU_i|-d_j(\mC_F)+1)^+\ge 1$. Therefore $k_{I_i}-j+1-\bigl(|\mU_i|-d_j(\mC_F)+1\bigr)^+
\le
k_{I_i}-j$.

Combining the inequalities above, the term with $\ell=j+1$ is at least as large as all other terms in the maximum from \cref{prop:readlower}, and its value is $k_{I_i}-j$. Hence
$|\mR_i|\ge k_{I_i}-j$, as claimed.
    \end{proof}



\begin{remark}\label{remark:bound-MDS}
We note that this bound coincides with the bound for stable linear MDS codes from~\cite{maturana2022convertible}. Stable codes attain the bound on the number of unchanged symbols, and in our analysis we do not distinguish the extreme case in which the default approach is optimal. 
\end{remark}

\begin{remark}\label{remark:readcost-LRC}
We now compare the read-cost bounds above with the bound for optimal locally repairable codes from~\cite{ge2025locally}, also independently obtained in~\cite{shi2025bounds}. For optimal locally repairable codes, lower bounds on the generalized Hamming weights are known~\cite{hao2020generalized}. Hence one can combine these lower bounds with Theorem~\ref{thm:readlowergen}. As before, we consider the symmetric setting
$k_{I_i}=k \quad \text{for all } i\in[\lambda], \qquad
    k_F=\lambda k, \qquad
    k_{-i}=(\lambda-1)k, \qquad
    n_F=n$, and we assume that all codes have locality $r$. Since the final code is optimal with locality, its minimum distance is
\begin{align*}
    d=n-\lambda k-\left\lceil \frac{\lambda k}{r}\right\rceil+2,
\end{align*}
and its generalized Hamming weights satisfy
\begin{align*}
    d_j(\mC_F)\ge d+j-1+\left\lfloor \frac{j-1}{r}\right\rfloor .
\end{align*}
Moreover, the bound on unchanged symbols gives
\begin{align*}
    |\mU_i|\le n-d-(\lambda-1)k+1
    -\left(\left\lceil\frac{(\lambda-1)k}{r}\right\rceil-1\right).
\end{align*}
Set
\begin{align*}
    \Delta
    :=n-2d-(\lambda-1)k+2
    -\left(\left\lceil\frac{(\lambda-1)k}{r}\right\rceil-1\right),
\end{align*}
and let $j^*\in[k]$ be the largest index such that $d_{j^*}(\mC_F)<|\mU_i|+1.$
Using the preceding upper bound on $|\mU_i|$, we get
\begin{align*}
    j^*+\left\lfloor\frac{j^*-1}{r}\right\rfloor
    \le |\mU_i|-d+1
    \le \Delta .
\end{align*}
If $\Delta\le 0$, then no such index $j^*$ exists, and Theorem~\ref{thm:readlowergen} gives $|\mR_i|\ge k.$
Assume now that $\Delta>0$. We first note that $j^*<k$. Indeed, if $j^*=k$, then the optimal locality parameters would imply
\begin{align*}
    n-\left\lceil
        \left\lceil \frac{\lambda k}{r}\right\rceil
        -\frac{k-1}{r}
    \right\rceil+1
    <
    \left\lceil
        \left\lceil \frac{\lambda k}{r}\right\rceil
        -\frac{(\lambda-1)k-1}{r}
    \right\rceil .
\end{align*}
Writing
$A:=\lceil \frac{\lambda k}{r}\rceil$,
this gives
\begin{align*}
    n
    \le
    2A
    -\left\lfloor \frac{k-1}{r}\right\rfloor
    -\left\lfloor \frac{(\lambda-1)k-1}{r}\right\rfloor
    -2 .
\end{align*}
Now write $k=qr+s$, where $0\le s<r$.
If $s=0$, then
\begin{align*}
    &A=\lambda q, \qquad
    \left\lfloor \frac{k-1}{r}\right\rfloor=q-1,\\
    &\left\lfloor \frac{(\lambda-1)k-1}{r}\right\rfloor=(\lambda-1)q-1.
\end{align*}
Thus
\begin{align*}
    n
    &\le
    2A
    -\left\lfloor \frac{k-1}{r}\right\rfloor
    -\left\lfloor \frac{(\lambda-1)k-1}{r}\right\rfloor
    -2\\ 
    &=2\lambda q-(q-1)-((\lambda-1)q-1)-2 \\
    &=\lambda q
    \le \lambda qr
    =\lambda k .
\end{align*}
If $s\ne 0$, then
\begin{align*}
    A=\lambda q+\left\lceil \frac{\lambda s}{r}\right\rceil,
    \qquad
    \left\lfloor \frac{k-1}{r}\right\rfloor=q,
\end{align*}
and
\begin{align*}
    \left\lfloor \frac{(\lambda-1)k-1}{r}\right\rfloor
    \ge
    (\lambda-1)q+\left\lfloor \frac{\lambda s}{r}\right\rfloor-1.
\end{align*}
Hence
\begin{align*}
    n
    &\le
    2A
    -\left\lfloor \frac{k-1}{r}\right\rfloor
    -\left\lfloor \frac{(\lambda-1)k-1}{r}\right\rfloor
    -2 \\
    &\le
    2\lambda q
    +2\left\lceil \frac{\lambda s}{r}\right\rceil
    -q-(\lambda-1)q
    -\left\lfloor \frac{\lambda s}{r}\right\rfloor
    +1-2 \\
    &=
    \lambda q
    +2\left\lceil \frac{\lambda s}{r}\right\rceil
    -\left\lfloor \frac{\lambda s}{r}\right\rfloor
    -1 \\
    &\le
    \lambda q+\lambda s
    \le
    \lambda(qr+s)
    =\lambda k .
\end{align*}
In both cases we obtain $n\le \lambda k$, contradicting the assumption $n>\lambda k$. Therefore $j^*<k$.

Let $a:=j^*-1$. Then
\begin{align*}
    a+\left\lfloor\frac{a}{r}\right\rfloor<\Delta .
\end{align*}
Let $a^*$ be the smallest integer satisfying
\begin{align*}
    a^*+\left\lfloor\frac{a^*}{r}\right\rfloor\ge \Delta .
\end{align*}
Then $j^*\le a^*$. Writing
\begin{align*}
    \Delta=(r+1)q+b, \qquad 0\le b<r+1,
    \qquad q=\left\lfloor\frac{\Delta}{r+1}\right\rfloor,
\end{align*}
we obtain
\begin{align*}
    a^*=rq+b
    =\Delta-\left\lfloor\frac{\Delta}{r+1}\right\rfloor .
\end{align*}
Consequently,
\begin{align*}
    |\mR_i|
    \ge k-j^*
    \ge k-a^*
    = k-\Delta+\left\lfloor\frac{\Delta}{r+1}\right\rfloor .
\end{align*}
Thus Theorem~\ref{thm:readlowergen} gives
\begin{align*}
|\mR_i|
\ge
\begin{cases}
k, 
& \Delta\le 0,\\[0.5em]
k-\Delta+\left\lfloor \dfrac{\Delta}{r+1}\right\rfloor,
& \Delta>0 .
\end{cases}
\end{align*}
This coincides with the bound from~\cite{ge2025locally}, also independently obtained in~\cite{shi2025bounds}, in this special setting. The coincidence relies on the optimal locally repairable code case. For other code families, especially when the generalized Hamming weights are larger, the bound from Theorem~\ref{thm:readlowergen} can be stronger. At the same time, it applies to arbitrary linear codes, including conversions between different code classes, which were not considered before.
% Let us compare these bounds on read cost for the special case of optimal
% locally repairable codes with bound from~\cite{ge2025locally}, which was independently obtained in~\cite{shi2025bounds}. For such codes, we know lower bounds on the
% generalized Hamming weights~\cite{hao2020generalized}, and therefore can substitute the closed-form
% bound on the number of unchanged symbols into the bound from Theorem~\ref{thm:readlowergen}. As before,  without loss of generality, we consider the case where
% $k_{I_i}=k$ for all $i\in[\lambda]$, $k_F=\lambda k$,
% $k_{-i}=(\lambda-1)k$, $n_F=n$, and the locality is $r$ for all codes. We use the fact that the final
% code is an optimal code with locality. Hence
% $d=n-\lambda k-\left\lceil \frac{\lambda k}{r}\right\rceil+2$ and, therefore, $d_j\geq d+j-1+\left\lfloor \frac{j-1}{r}\right\rfloor$.  As a result, we get $|\mU_i|\leq n-d-(\lambda-1)k+1-(\lceil\frac{(\lambda-1)k}{r}\rceil-1)$. 

% Let us denote $\Delta=n-2d-(\lambda-1)k+2-(\lceil\frac{(\lambda-1)k}{r}\rceil-1)$ and let $j^*\in[k]$ be the largest index for which $d_{j^*}<|\mU_i|+1$. It is clear that $j^{*}+\lfloor\frac{j^{*}-1}{r}\rfloor\leq|\mU_i|-d+1\leq\Delta$. 

% When $\Delta\leq0$ no such $j^{*}$ exist and, according to Theorem~\ref{thm:readlowergen}, we have $|\mR_i|\geq k$. 

% Now assume that $\Delta>0$. We focus on the case when $j^{*}<k$ as $j^*=k$ for optimal codes with locality  would imply $n-\left\lceil
% \left\lceil \frac{\lambda k}{r}\right\rceil
% -\frac{k-1}{r}
% \right\rceil
% +1
% <
% \left\lceil
% \left\lceil \frac{\lambda k}{r}\right\rceil
% -\frac{(\lambda-1)k-1}{r}
% \right\rceil$. Denoting $A=\left\lceil \frac{\lambda k}{r}\right\rceil$, this would give $n
% \le
% 2A
% -
% \left\lfloor \frac{k-1}{r}\right\rfloor
% -
% \left\lfloor \frac{(\lambda-1)k-1}{r}\right\rfloor
% -2$. Write $k=qr+s$, $0\le s<r$.

% If \(s=0\), then \(A=\lambda q\), $\left\lfloor \frac{k-1}{r}\right\rfloor=q-1$, $\left\lfloor \frac{(\lambda-1)k-1}{r}\right\rfloor
% =(\lambda-1)q-1$. As a result, we have $n\leq 2A
% -
% \left\lfloor \frac{k-1}{r}\right\rfloor
% -
% \left\lfloor \frac{(\lambda-1)k-1}{r}\right\rfloor
% -2
% =
% 2\lambda q-(q-1)-((\lambda-1)q-1)-2
% =
% \lambda q
% \le \lambda qr
% =
% \lambda k$.


% If \(s\ne 0\), then $A=\lambda q+\left\lceil \frac{\lambda s}{r}\right\rceil$, $\left\lfloor \frac{k-1}{r}\right\rfloor=q$. Moreover, it can be easily shown that $\left\lfloor \frac{(\lambda-1)k-1}{r}\right\rfloor
% \geq
% (\lambda-1)q+\left\lfloor \frac{\lambda s}{r}\right\rfloor-1$. As a result, we get $2A
% -
% \left\lfloor \frac{k-1}{r}\right\rfloor
% -
% \left\lfloor \frac{(\lambda-1)k-1}{r}\right\rfloor
% -2 \le
% 2\lambda q
% +
% 2\left\lceil \frac{\lambda s}{r}\right\rceil
% -q
% -(\lambda-1)q
% -
% \left\lfloor \frac{\lambda s}{r}\right\rfloor
% +1
% -2                                                
% =
% \lambda q
% +
% 2\left\lceil \frac{\lambda s}{r}\right\rceil
% -
% \left\lfloor \frac{\lambda s}{r}\right\rfloor
% -1                                                
% \le
% \lambda q+\lambda s
% \le
% \lambda(qr+s)
% =
% \lambda k$.

% Thus in all cases we get \(n\le \lambda k\), which contradicts the
% assumption \(n>\lambda k\). Hence \(j^*<k\).

% Let us introduce $a=j^*-1$. Clearly, $a+\lfloor\frac{a}{r}\rfloor<\Delta$. Let $a^*$ be the smallest integer for that $a^*+\left\lfloor \frac{a^*}{r}\right\rfloor \ge \Delta$ and, as a result, $j^*\leq a^*$. Writing $\Delta=(r+1)q+b$, where $0\le b < r+1$ and $q=\left\lfloor \frac{\Delta}{r+1}\right\rfloor$, we get $a^*=rq+b
% =
% \Delta-\left\lfloor \frac{\Delta}{r+1}\right\rfloor$.
% Consequently, $|\mR_i|
% \ge
% k-j^*
% \ge
% k-a^*
% =
% k-\Delta+\left\lfloor \frac{\Delta}{r+1}\right\rfloor$.

% Hence, the bound from Theorem~\ref{thm:readlowergen} has the following shape.
% \[
% |\mR_i|
% \ge
% \begin{cases}
% k, 
% & \Delta\le 0,\\[0.5em]
% k-\Delta+\left\lfloor \dfrac{\Delta}{r+1}\right\rfloor,
% & \Delta>0 .
% \end{cases}
% \]

% This coincides with the bound from~\cite{ge2025locally}, which was also independently obtained in~\cite{shi2025bounds}, in this special setting. However, this coincidence relies on the final code being an optimal code with locality and on the corresponding generalized-Hamming-weight bound being tight. When this bound is not tight, or when the codes with locality are not optimal, our bound can be stronger.

% An example of the latter is given by the conversion from two binary codes $(a,a,b,b)$ with locality one and two information symbols into one binary Hamming code of length seven. The initial codes clearly have locality one, while the final code has locality three. Theorem~\ref{thm:readlowergen} gives $(|\mR_i|\geq 2)$, showing that only the trivial read cost is possible. In contrast, the bounds from~\cite{ge2025locally} and~\cite{shi2025bounds} give only $|\mR_i|\geq 1$, leaving open the possibility of a non-trivial conversion, which is ruled out by our bound. 

% We also note that our bound applies to arbitrary linear codes, including conversions between different code classes, which were not considered previously.
\end{remark}
As in Subsection~\ref{subsec:boundsunchanged} we can reduce the bound from Proposition~\ref{prop:readlower} to a bound only taking into account the minimum distance of the final code, instead of going through all the generalized Hamming weights.

\begin{corollary}\label{cor:readloweru}
Let $(\mC_{\textbf{I}},\mC_F,\sigma)$ be a convertible code. Fix $i\in [\lambda]$, and let $\mR_i$ and $\mU_i$ be the sets of read and unchanged symbols corresponding to $\mC_{I_i}$, respectively. Then
\begin{align*} |\mR_i|\ge k_{I_i}-\bigl(|\mU_i|-d_F+1\bigr)^+. \end{align*}
\end{corollary}

\begin{proof}
This is the special case $j=1$ of Proposition~\ref{prop:readlower}, using $d_1(\mC_F)=d(\mC_F)=d_F$.
\end{proof}

\begin{remark}\label{rem:read-ghw-improvement}
Fix $i\in [\lambda]$ and set $U_i:=|\mU_i|$. For $j\in [k_{I_i}]$ we define
\begin{align*}
    B_j(U_i):=k_{I_i}-j+1-\bigl(U_i-d_j(\mC_F)+1\bigr)^+.
\end{align*}
Then Proposition~\ref{prop:readlower} states that $|\mR_i|\ge \max_{1\le j\le k_{I_i}} B_j(U_i)$, while Corollary~\ref{cor:readloweru} is exactly the special case $|\mR_i|\ge B_1(U_i)$. As a result, the bound taking into account generalized Hamming weights is never weaker than the bound based only on the minimum distance. For a fixed $j\ge 2$, it is strictly stronger whenever
\begin{align*}
    d_j(\mC_F)>d_F+j-1 \quad \text{and} \quad U_i\ge d_F+j-1.
\end{align*}
The first condition means that the generalized Hamming weight hierarchy has a non-trivial jump at level $j$, while the second ensures that we have enough unchanged symbols for this extra information to affect the term $(U_i-d_j(\mC_F)+1)^+$. Thus a jump in the generalized Hamming weights is necessary for an improvement over the bound with only the minimum distance, but it only becomes visible in the read-cost estimate once $|\mU_i|$ is sufficiently large.
\end{remark}




Suppose we do not know the number of unchanged symbols of our code conversion. Then we can combine Corollary~\ref{cor:unchanged1} with Corollary~\ref{cor:readloweru} to obtain the following result.


\begin{corollary}\label{cor:readlowergen}
Let $(\mC_{\textbf{I}},\mC_F,\sigma)$ be a convertible code, and let $i\in [\lambda]$. Then
\begin{align*}
|\mR_i|\ge
\begin{cases}
k_{I_i}, & \text{if } n_F-2d_F-\sum_{j\ne i} k_{I_j}+2\le 0,\\[2pt]
k_{I_i}-(n_F-2d_F\\-\sum_{j\ne i} k_{I_j}+2), & \text{if } n_F-2d_F-\sum_{j\ne i} k_{I_j}+2>0.
\end{cases}
\end{align*}
\end{corollary}

\begin{proof}
By Corollary~\ref{cor:unchanged1},
\begin{align*} |\mU_i|\le n_F-d_F-\sum_{j\ne i} k_{I_j}+1. \end{align*}
Substituting this bound into Corollary~\ref{cor:readloweru} gives
\begin{align*} |\mR_i|\ge k_{I_i}-\bigl(n_F-2d_F-\sum_{j\ne i} k_{I_j}+2\bigr)^+. \end{align*}
Writing this bound piecewise gives the claimed inequality.
\end{proof}

\begin{remark}
We note that this bound coincides with the bound for arbitrary linear MDS convertible codes from~\cite{maturana2022convertible}. As in Remark~\ref{remark:bound-MDS}, we do not distinguish the extreme case in which the default approach is optimal. 
\end{remark}


\section{Conversion of Reed--Muller Codes}\label{sec:RM}

\subsection{Definitions and Properties}
We start by giving preliminary definitions and basic properties of \emph{Reed--Muller} codes; see, for example,~\cite{macwilliams1977theory, abbe2020reed}. In what follows, for integers $r \ge 0$ and $m \ge 1$, we denote by $\F_2[X_1,\dots,X_m]_{\le r}^\times$
the vector space of square-free polynomials in $X_1,\dots,X_m$ of total degree at most $r$. 
We also let $\mathcal{P}(m)$ denote the list of vectors in $\F_2^m$ sorted in lexicographic order. For example,
for $m=3$ we have
\begin{align} \label{ptsRM}
\mathcal{P}(3)=
\bigl(&(0,0,0),(0,0,1),(0,1,0),(0,1,1),
\\ &(1,0,0),(1,0,1),(1,1,0),(1,1,1)\bigr) \nonumber
\end{align}
\begin{definition}\label{def:rm}
Let $r \ge 0$ and $m \ge 1$ be integers. Let $n=2^m$ and $\mP(m)=(a_1,...,a_n)$.
The \textbf{Reed--Muller code} with parameters $(r,m)$ is
$\RM(r,m):=\{(p(a_1),...,p(a_n)) : p \in \F_2[X_1,...,X_m]_{\le r}^\times\}$.
\end{definition}

Some of the properties of Reed--Muller codes that we will need in the sequel are summarized as follows~\cite{abbe2020reed}.

\begin{theorem} \label{thm:RMdim}
Let $r \ge 0$ and $m \ge 1$ be integers, and $n=2^m$. 
\begin{itemize}
    \item[(i)] $\RM(r,m)$ is a linear code of dimension $\sum_{i=0}^r \binom{m}{i}$;
    \item[(ii)] The minimum distance of $\RM(r,m)$ is
$2^{m-r}$;
    \item[(iii)] For $r\ge 1$ and $m \ge 1$ we have $\RM(r,m) = \RM(r,m-1) \ps \RM(r-1,m-1),$ where for subspaces $\mC, \mD \le \F_q^n$ we let $$\mC \ps \mD:=\{(x,x+y) : x \in \mC, \, y \in \mD \} \le \F_q^{2n}$$ be the \textbf{Plotkin sum};
    \item[(iv)] The dual of $\RM(r,m)$ is $\RM(r,m)^\perp = \RM(m-r-1,m)$.
\end{itemize}
\end{theorem}

%The reason why we are interested in Reed--Muller codes in the context of convertible codes is because of their good \textit{locality} with \textit{availability} properties. We do not go into detail here, but refer the reader to~\cite{}.


\subsection{Reed--Muller Convertible Codes}

In this subsection, we explain how to construct Reed--Muller convertible codes for $\lambda=2$. This construction can be extended recursively for larger $\lambda$.

Let $r \ge 1$ and $m \ge 1$, and denote by $G_{\RM(r,m)}$ a generator matrix of $\RM(r,m)$.

\begin{example}
Let $r=2$ and $m=3$. Then $n=8$ and $\RM(2,3)$ is generated by the evaluations 
of the polynomials in $\{1,X_1, X_2, X_3,X_1X_2,X_1X_3,X_2X_3\}$ at the points listed in~\eqref{ptsRM}. For a multivariate polynomial $p$, denote by $\textup{ev}(p)$ the vector $(p(a_1),\dots,p(a_n)) \in \F_2^n$.
Therefore a generator matrix $G_{\RM(2,3)}$ of $\RM(2,3)$ is
\begin{small}
\begin{align} \label{eq:RMex}
\begin{bmatrix}
        \textup{ev}(1) \\
        \textup{ev}(X_1) \\
        \textup{ev}(X_2) \\
        \textup{ev}(X_3) \\
        \textup{ev}(X_1X_2) \\
        \textup{ev}(X_1X_3) \\
        \textup{ev}(X_2X_3) \\
    \end{bmatrix}
    = 
    \begin{bmatrix}
1  & 1 & 1 & 1 & 1 & 1 & 1 & 1 \\
0  & 0 & 0 & 0 & 1 & 1 & 1 & 1 \\
0  & 0 & 1 & 1 & 0 & 0 & 1 & 1 \\
0  & 1 & 0 & 1 & 0 & 1 & 0 & 1 \\
0  & 0 & 0 & 0 & 0 & 0 & 1 & 1 \\
0  & 0 & 0 & 0 & 0 & 1 & 0 & 1 \\
0  & 0 & 0 & 1 & 0 & 0 & 0 & 1
\end{bmatrix} 
= \begin{bmatrix}
        G_{\RM(1,3)} \\
        \textup{ev}(X_1X_2) \\
        \textup{ev}(X_1X_3) \\
        \textup{ev}(X_2X_3) \\
    \end{bmatrix}.
\end{align}
\end{small}
\end{example}


\begin{remark} \label{rem:rmsplit}
From Theorem~\ref{thm:RMdim} (iii) we know that $\RM(r,m) = \RM(r,m-1) \ps \RM(r-1,m-1)$. In terms of generator matrices, this means there exists a generator matrix of the form
\begin{align*}
    G_{\RM(r,m)} = \begin{pmatrix}
        G_{\RM(r,m-1)} & G_{\RM(r,m-1)} \\
        0 & G_{\RM(r-1,m-1)}
    \end{pmatrix}.
\end{align*}

Moreover,
\begin{align*}
    G_{\RM(r,m-1)} = \begin{pmatrix}
        G_{\RM(r-1,m-1)} \\
        A
    \end{pmatrix},
\end{align*}
where $A$ is a matrix formed by evaluating the monomials $p \in \F_2[X_2,...,X_m]_{ \le r}^\times$ of degree~$r$ at the points in $\mP(m-1)$. In particular, $A$ has $\sum_{i=0}^{r-1} \binom{m-1}{i}$ zero columns (because any vector in $\mP(m-1)$ of weight at most $r-1$ will be 0 when evaluated at a monomial of degree $r$). Therefore, by performing row operations on $G_{\RM(r,m)}$ we obtain a matrix of the form
\begin{align} \label{eq:rmm}
    \Tilde{G}_{\RM(r,m)}= \begin{pmatrix}
        G_{\RM(r,m-1)} & 0 \\
         & A \\
       0 & G_{\RM(r-1,m-1)}
    \end{pmatrix}
\end{align}
which is clearly also a generator matrix of $\RM(r,m)$.
\end{remark}

%We are now ready to introduce Reed--Muller convertible codes more generally.

\begin{theorem} \label{thm:rmc}
    Let $r \ge 1$ and $m \ge 2$. Let $\mC_{I_1}=\RM(r,m-1)$, $\mC_{I_2}=\RM(r-1,m-1)$, $\mC_{\textbf{I}}=\mC_{I_1} \times \mC_{I_2}$, and $\mC_F=\RM(r,m)$. There exists a convertible code $(\mC_{\textbf{I}},\mC_F,\sigma)$ with 
    \begin{align*}
        &|\mU_1| = n_{I_1}, \quad |\mU_2| = k_{I_2}, \\
        &|\mR_1| \le  k_{I_1}, \quad |\mR_2| = \min\{k_{I_2},n_{I_2}-k_{I_2}\}.
    \end{align*}
\end{theorem}
\begin{proof}
    We use the form of the generator matrix of $\RM(r,m)$ as in~\eqref{eq:rmm} adopting also the notation used in Remark~\ref{rem:rmsplit}. The conversion procedure is defined by copying the coordinates corresponding to the first block $\RM(r,m-1)$, copying the $k_{I_2}$ coordinates of $\RM(r-1,m-1)$ corresponding to the zero columns of $A$, and computing the remaining coordinates using the displayed generator-matrix representation. More precisely, the remaining coordinates are obtained from the nonzero columns of $A$ and the corresponding coordinates of $\RM(r-1,m-1)$; hence their computation requires reading only the coordinates of $\mathcal{C}_{I_2}$ not copied unchanged, together with at most $k_{I_1}$ coordinates from $\mathcal{C}_{I_1}$. Note that we can leave all symbols from $\mC_{I_1}$ unchanged, i.e., $|\mU_1|  = n_{I_1}$. Now $A$ has $\sum_{i=0}^{r-1} \binom{m-1}{i}=k_{I_2}$ zero columns. Thus we can leave the $k_{I_2}$ coordinates corresponding to these zero columns unchanged, and we have to read $\min\{k_{I_2},n_{I_2}-k_{I_2}\}$ symbols in order to generate the remaining $n_{I_2}-k_{I_2}$ columns. In order to generate the  $n_{I_2}-k_{I_2}$ nonzero columns of $A$, we need to read at most $k_{I_1}$ symbols from $\mC_{I_1}$. 
\end{proof}
%The Reed--Muller convertible code from Theorem~\ref{thm:rmc} is access optimal with respect to the general bounds we derived in Section~\ref{sec:gen}.

\subsection{Comparison with the obtained bounds}\label{sec:comp}

In this section, we consider the convertible code from Theorem~\ref{thm:rmc} for $m = r + 2 \ge 4$, and compare its read and write costs to the bounds in Subsections~\ref{subsec:boundsunchanged} and~\ref{subsec:boundsread}. We focus on the case $m=r+2$ because then $\mC_F^\perp=\RM(1,m)$, whose generalized
Hamming weights are explicit. This allows a closed-form comparison with the bounds from
Section~\ref{sec:gen}. To this end, we state the corresponding bounds from Corollaries~\ref{cor:unchanged1},~\ref{cor:unchanged2},~\ref{cor:readloweru}, and \cref{thm:readlowergen}, and derive the conditions under which these bounds are met with equality.
\begin{align*}
        &n_{I_1}=n_{I_2}=2^{m-1}, \quad k_{I_1} = \sum_{i=0}^{r} \binom{m-1}{i}=2^{m-1}-1,\\
        &k_{I_2} = \sum_{i=0}^{r-1} \binom{m-1}{i} = 2^{m-1}-m, \quad n_{F}=2^{m},\\
        &k_{F} = \sum_{i=0}^{r} \binom{m}{i},\quad d_F=2^{m-r}=2^2, \quad d_F^\perp = 2^{r+1}=2^{m-1}.
    \end{align*}
    The bound of Corollary~\ref{cor:unchanged1} evaluated at the above parameters and for $m=r+2$ reads as follows:
\[
\scalebox{0.85}{$
\aligned
|\mathcal{U}_1|
&\le 2^m - 2^2 - \sum_{i=0}^{m-3} \binom{m-1}{i} + 1
= 2^{m-1} + m - 3
\Longrightarrow |\mathcal{U}_1| \le n_{I_1}, \\
|\mathcal{U}_2|
&\le 2^m - 2^2 - \sum_{i=0}^{m-2} \binom{m-1}{i} + 1
= 2^{m-1} - 2
\Longrightarrow |\mathcal{U}_2| \le n_{I_2} - 2 .
\endaligned
$}
\]
Now note that $d_F^\perp = 2^{m-1} \le (2^{m-1}-1)+1 = k_{I_{1}}+1,$ so Corollary~\ref{cor:unchanged2} is not applicable to $\mC_{I_1}$. For $\mC_{I_2}$ we have $$d_F^\perp = 2^{m-1} > (2^{m-1}-m)+1 = k_{I_{2}}+1,$$ so Corollary~\ref{cor:unchanged2} applies, and we have $|\mU_2| \le k_{I_2}$. This shows that the Reed--Muller convertible code meets the bounds on the number of unchanged symbols with equality, and thus our conversion has optimal write cost.

We now compare these write-cost bounds with the generalized-Hamming-weight bound from Proposition~\ref{prop:ghw-unchanged-rewrite}. Since $\mC_F=\RM(m-2,m)$, we have $\mC_F^\perp=\RM(1,m)$. The generalized Hamming weights of $\RM(1,m)$ are given by
\begin{align*}
    d_s(\RM(1,m))=2^m-2^{m-s}, \qquad 1\le s\le m,
\end{align*}
and $d_{m+1}(\RM(1,m))=2^m$. By Wei duality (\cref{thm:wei-duality-rewrite}), the generalized Hamming weights of $\mC_F$ are therefore the integers in $$[2^m] \setminus (\{2^{m-s} +1\, : 1 \leq s \leq m\} \cup \{1\}).$$
% excluding
% \begin{align*}
%     1,2,3,5,9,\dots,2^{m-1}+1.
% \end{align*}
In particular,
\begin{align*}
    d_{k_{I_2}}(\mC_F)=2^{m-1} \quad \text{and} \quad d_{k_{I_1}}(\mC_F)=2^{m-1}+m.
\end{align*}
Hence Proposition~\ref{prop:ghw-unchanged-rewrite} yields
\begin{align*}
    |\mU_1|\le 2^m-2^{m-1}=2^{m-1}=n_{I_1}
\end{align*}
and
\begin{align*}
    |\mU_2|\le 2^m-(2^{m-1}+m)=2^{m-1}-m=k_{I_2}.
\end{align*}
Thus the generalized-Hamming-weight bound does not improve the comparison for $|\mU_1|$, since the result is again truncated by the trivial bound $|\mU_1|\le n_{I_1}$, but it does sharpen the distance-only bound for $|\mU_2|$ from $2^{m-1}-2$ to $2^{m-1}-m$. This coincides with the bound obtained above from Corollary~\ref{cor:unchanged2}.

For the read cost, we first evaluate the minimum-distance bound from Corollary~\ref{cor:readloweru}. We have $\delta_1 = n_{I_1}-d_F+1=2^{m-1}-2^2+1 \ge 0$ and so $|\mR_1| \ge k_{I_1}-2^{m-1}+2^2-1=2$. Thus, for $|\mR_1|$, this bound is not met with equality by our construction.

We have $\delta_2 = k_{I_2}-d_F+1=2^{m-1}-m-2^2+1 > 0$ for $m \ge 4$ and so $|\mR_2| \ge 3$. Since $\min\{n_{I_2}-k_{I_2},k_{I_2}\}=\min\{m,2^{m-1}-m\}=m$, this bound is not met with equality for all $m\ge 4$.

The generalized-Hamming-weight bound from \cref{thm:readlowergen} gives a better comparison. Since $|\mU_1|=2^{m-1}$ and $d_{k_{I_2}}(\mC_F)=2^{m-1}$ we have  $d_{k_{I_2}}(\mC_F) < |\mU_1| +1$. Furthermore, since $k_{I_2} < k_{I_1}$ and from the strict monotonicity of the generalized Hamming weight, $|\mU_1| + 1 \leq d_{k_{I_2}+1}(\mC_F)$. By \cref{thm:readlowergen} we therefore get
$|\mR_1| \geq k_{I_1} - k_{I_2} = 2^{m-1} - 1 - (2^{m-1} -m) = m-1$.


Thus the generalized-Hamming-weight bound improves the lower bound on $|\mR_1|$ from $2$ to $m-1$. This is a substantial improvement, but our construction still has $|\mR_1|\le k_{I_1}=2^{m-1}-1$, so the gap remains large.

For $|\mR_2|$, the improvement is stronger. If $m\ge 5$, then among the integers $1,\dots,2^{m-1}-m$, exactly the elements of
\begin{align*}
    \{1\}\cup \{2^{m-s}+1 : 2\le s\le m\}
\end{align*}
are excluded from the weight hierarchy of $\mC_F$. Hence
\begin{align*}
    d_{2^{m-1}-2m}(\mC_F)=2^{m-1}-m = k_{I_2} = |\mU_2|.
\end{align*}
Therefore, letting $j :=2^{m-1}-2m=k_{I_2}-m$, by strict monotonicity of the generalized Hamming weight and the fact that $j < k_{I_2}$, we get from  \cref{thm:readlowergen} that
$$|\mR_2| \geq k_{I_2} - j = m$$

Since our construction satisfies $|\mR_2|=\min\{k_{I_2},n_{I_2}-k_{I_2}\}=m$, the generalized-Hamming-weight bound is sharp for $|\mR_2|$ whenever $m\ge 5$. In the remaining case $m=4$, a direct evaluation of \cref{thm:readlowergen} yields only $|\mR_2|\ge 3$.

The comparisons above are summarized in the following two tables. The first table collects the bounds on the number of unchanged symbols. In the last row we recall the number of unchanged symbols attained by the Reed--Muller conversion from Theorem~\ref{thm:rmc}.
\begin{table}[h!]\centering
\caption{Comparison of bounds on unchanged symbols for the Reed--Muller conversion with $m=r+2\ge 4$.}
\label{tab:rm-unchanged-comparison}
\begin{scalebox}{0.8}{
\begin{tabular}{l|l|l}
\toprule
Result used & $|\mU_1|$ & $|\mU_2|$ \\
\midrule
Corollary~\ref{cor:unchanged1} & $\le n_{I_1}$ & $\le n_{I_2}-2$ \\
Corollary~\ref{cor:unchanged2} & not applicable & $\le k_{I_2}=2^{m-1}-m$ \\
Proposition~\ref{prop:ghw-unchanged-rewrite} & $\le n_{I_1}=2^{m-1}$ & $\le k_{I_2}=2^{m-1}-m$ \\
\midrule
Construction & $n_{I_1}=2^{m-1}$ & $k_{I_2}=2^{m-1}-m$ \\
\bottomrule
\end{tabular}
}\end{scalebox}

\end{table}

The second table collects the corresponding lower bounds on read symbols. Here the last row gives the read cost of the construction from Theorem~\ref{thm:rmc}.
\begin{table}[h!]
\centering
\caption{Comparison of bounds on read symbols for the Reed--Muller conversion with $m=r+2\ge 4$.}
\label{tab:rm-read-comparison}
\begin{scalebox}{0.8}{
\begin{tabular}{l|l|l}
\toprule
Result used & $|\mR_1|$ & $|\mR_2|$ \\
\midrule
Corollary~\ref{cor:readloweru} & $\ge 2$ & $\ge 3$ \\
Corollary~\ref{cor:readlowergen} & $\ge \max\{0,5-m\}$ & $\ge \max\{0,5-m\}$ \\
 Theorem~\ref{thm:readlowergen} & $\ge m-1$ & $\ge m$ for $m\ge 5$; $\ge 3$ for $m=4$ \\
\midrule
Construction & $\le k_{I_1}=2^{m-1}-1$ & $m$ \\
\bottomrule
\end{tabular}
}\end{scalebox}
\end{table}


\begin{remark}
The above example shows that for the case where $m = r+2 \ge 4$, our construction of Reed--Muller convertible codes has optimal write cost: the bounds for $|\mU_1|$ and $|\mU_2|$ are met with equality. The generalized-Hamming-weight refinement does not change the comparison for $|\mU_1|$, but it improves the minimum distance-only bound for $|\mU_2|$ from $2^{m-1}-2$ to $2^{m-1}-m=k_{I_2}$, which is again attained by our construction. For the read cost, the generalized-Hamming-weight bounds improve the comparison more substantially: they raise the lower bound on $|\mR_1|$ from $2$ to $m-1$, and they give the sharp bound $|\mR_2|\ge m$ for all $m\ge 5$. Thus generalized Hamming weights capture the behavior of $|\mR_2|$, but the gap for $|\mR_1|$ remains large.
    % For $|\mR_1|$, even though our construction does not meet the lower bound presented in Proposition~\ref{prop:readlower}, one cannot do better than $|\mR_1| = k_{I_1}$ here. The reason for this is that we have $|\mU_1|=n_{I_1}$, $|\mU_2|=k_{I_2}$ and $d_F^\perp=d_{I_1}^\perp=2^{m-1}$. Note that the $k_{I_2}$ columns of $G_{\RM(r-1,m-1)}$ for which $A$ is zero (so $\mU_2$) correspond to the Reed--Muller code $\RM(r-1,r-1)$, and thus consists of linearly independent columns. Moreover, any set of $2^{m-1}-1$ symbols from $\mU_1$ are linearly independent (since $d_{I_1}^{\perp} = 2^{m-1}$), and also any set of $2^{m-1}-1$ symbols from $\mC_F$ are linearly independent (since $d_F^{\perp} = 2^{m-1}$). Therefore it is not possible to generate the new symbols in $\mC_F$ from fewer than $k_{I_1}=2^{m-1}-1$ symbols from $\mC_{I_1}$, as this would give a dual codeword in $\mC_F^\perp$ of weight less than~$2^{m-1}$.

    % For $|\mR_2|$ we also do not meet the bound, and we do not have an explanation of why it is not possible to meet it.
\end{remark}


\section{Discussion and Future Directions}\label{sec:concl}
In this paper, we studied scalar linear code conversion in the merge regime from a general perspective. We developed a linear-algebraic framework for bounding the read and write costs of conversion between arbitrary linear codes. The resulting bounds are universal in the sense that they do not rely on structural properties specific to MDS codes, locally repairable codes, or other particular code families. Instead, they provide general access-cost lower bounds that can be used as a baseline for evaluating code-conversion constructions. We further refined these bounds using generalized Hamming weights, which allows the estimates to capture finer structural properties of the final code. In particular, when the generalized Hamming weight hierarchy has nontrivial jumps, the refined bounds can improve upon bounds obtained from the minimum distance alone and recover known bounds in several special cases.

Motivated by the role of Reed--Muller codes in distributed storage systems with heterogeneous request patterns, we also introduced an explicit merge-conversion procedure for Reed--Muller codes. The construction exploits the recursive structure of Reed--Muller codes through the Plotkin construction. In the parameter regime analyzed in this paper, the construction attains the derived bounds on the write cost, showing that Reed--Muller codes can support efficient conversion while retaining their algebraic structure and their relevance to heterogeneous-request storage models. For the read cost, the generalized-Hamming-weight bounds give a sharp estimate for one part of the conversion, while a gap remains for the other part. Closing this gap, or proving that it is inherent, is a natural open problem.

Several further directions remain open. First, it would be interesting to develop universal lower bounds on the bandwidth cost of linear code conversion, possibly by extending the generalized-Hamming-weight approach to vector codes and partial downloads. Second, generalized Hamming weights may be useful for refining existing access-cost bounds for specific families of convertible codes, such as locally repairable codes, codes with availability, or other codes whose weight hierarchies are sufficiently well understood. Finally, while this paper focuses on linear conversion procedures, the study of nonlinear code conversion remains largely unexplored and may reveal conversion mechanisms that are not captured by the present linear framework.


\balance
\sloppy
\printbibliography


\end{document}
